\chapter{奖励受限下的强化学习}
\label{chap:rl-reward}

一台配送机器人能够移动和避障，我们还需要告诉它怎样才算完成一次好的配送。若只在签收时给分，机器人可能拿到终局成绩，却难以判断早先哪次核对出了错。若每移动一米都给分，它又可能靠绕路增加得分，而不把包裹交给收件人。请人示范路线可以提供一组可用动作，但单凭这些动作，仍难判断示范者看重的是速度、能耗还是包裹完好。这里分别缺少动作贡献的依据、完整的计分标准和足以确定目标的信息，不能用同一种奖励处理办法回答。

第~\ref{chap:rl-learning}章虽然不知道期望奖励，仍假定交互后能读到奖励样本；第~\ref{chap:rl-information}章检查状态信息和数据覆盖是否足以支持决策。本章保留这些学习工具，改变它们所接收的反馈。我们要分别判断：没有可信的标量奖励时，其他记录能约束哪些行为目标；只有轨迹总分时，怎样用它更新局部决策；训练分数提高时，实际行为是否满足原先的要求。图~\ref{fig:rl-reward-map}说明这三个问题如何共同支持反馈学习。它们是可以叠加的限制，而不是必须依次排除的三种故障。

\begin{figure*}[htbp]
  \centering
  % Knowledge dependencies (solid) and feedback iteration (dashed).
\begin{tikzpicture}[
  x=1mm,y=1mm,font=\figurefont\fontsize{8}{10}\selectfont,
  box/.style={draw=BookRule,fill=BookPaper,minimum height=11mm,
    text width=27mm,align=center,inner sep=2mm,line width=.7pt},
  flow/.style={-{Stealth[length=2mm]},draw=BookMuted,line width=.8pt},
  loop/.style={flow,dashed},
  lab/.style={fill=white,inner sep=1mm,font=\figurefont\fontsize{7}{9}\selectfont}
]
  \node[box,fill=FigureInput!10] (limits) at (16,28) {三类奖励限制};
  \node[box] (obs) at (64,51) {观测\\可推断什么};
  \node[box] (credit) at (64,28) {归因\\如何分配结果};
  \node[box] (goal) at (64,5) {目标匹配\\行为是否符合要求};
  \node[box,fill=FigureModel!10] (update) at (112,39) {整合与更新};
  \node[box,fill=FigureOutput!10] (check) at (112,5) {独立检验\\继续或停止};
  \node[box,text width=20mm] (query) at (153,39) {定向查询\\新反馈};
  \draw[flow] (limits.east) -- ++(5,0) |- (obs.west);
  \draw[flow] (limits) -- (credit);
  \draw[flow] (limits.east) -- ++(5,0) |- (goal.west);
  \draw[flow] (obs.east) -- ++(6,0) |- (update.west);
  \draw[flow] (credit.east) -- ++(6,0) |- (update.west);
  \draw[flow] (goal) -- (check);
  \draw[loop] (update) -- node[lab,right] {新行为} (check);
  \draw[loop] (check.east) -- (153,5)
    -- node[lab,right,align=left] {有可补缺口\\且有预算} (query.south);
  \draw[loop] (query) -- node[lab,above] {反馈} (update);
  \draw[loop] (check.south) -- ++(0,-10)
    node[below,align=center] {满足要求／预算耗尽／识别受限\\分别记录停止结论};
\end{tikzpicture}

  \caption{三类奖励限制与反馈学习的关系。实线表示观测、归因和目标匹配分别为反馈整合、策略更新及行为检验提供依据，不规定三者的解决顺序；虚线表示取得新行为和新反馈后进入下一轮。停止可能因为证据已满足要求，也可能因为预算耗尽或关键目标信息无法补足，三者支持的结论不同。}
  \label{fig:rl-reward-map}
\end{figure*}

本章仍用$r_t$表示执行$a_t$后的奖励，用$J(\pi)$表示未归一化的期望折扣回报。为观察反馈落在哪个环节，我们另设一个四步配送任务：机器人依次核对包裹、装载、搬运和交付，状态包含包裹身份、位置及完好情况。这个案例不沿用前文随机网格的精确动力学；除特别说明外，反馈在交付后才给出。本章用于计算的人工数字都是教学假设。

% Outline: 5.1
\section{奖励信号的三类限制}
\label{sec:rl-reward-limits}

% Outline: 5.1.1
\subsection{奖励观测受限}
\label{sec:rl-reward-observation}

机器人不知道一条路线的平均耗电量时，仍可在执行后读取电表，再用样本估计期望奖励。只要电表记录可信，这属于前文的未知奖励学习。若目标改为“把包裹妥善交给收件人”，日志却只记录移动距离，情况就不同了：距离既没有说明交给了谁，也没有说明包裹是否破损。此时缺少表达目标的观测，而不只是均值估计不够精确；增加同类距离记录仍不能判断交付质量。

\begin{definition}[替代反馈的形式与粒度]
  \label{def:rl-reward-feedback}
  \term{示范}（demonstration）记录决策者实际采取的行为，可以是某状态上的动作，也可以是整条轨迹。\term{评分}（rating）按照约定尺度给被评价对象一个数值。\term{偏好}（preference）记录在相同比较条件下，评价者如何排列两个或多个对象。\term{验证反馈}（verification feedback）判断对象是否满足指定条件，并可附带检查证据。反馈的\term{粒度}（granularity）是这次判断实际评价的单位，可以是动作、步骤、轨迹片段或完整轨迹。记录到达的时间不必等于被评价动作发生的时间。
\end{definition}

这些形式保留的信息不同。示范告诉我们“当时做了什么”，没有直接评价未选动作；评分保留数值差距，但比较差值需要统一尺度。偏好可以只给出顺序，不给绝对分数；验证器的通过标记则只支持被检查条件。表~\ref{tab:rl-reward-forms}还列出查询权限，因为读取已有示范与继续向专家提问是两种访问能力，前者不能代替后者。

\begin{table*}[htbp]
  \centering\small
  \begin{tabularx}{\linewidth}{@{}>{\raggedright\arraybackslash}p{17mm}*{5}{>{\raggedright\arraybackslash}X}@{}}
    \toprule
    形式 & 反馈对象 & 粒度 & 噪声来源 & 查询权限 & 可推断对象 \\
    \midrule
    示范 & 已执行动作与上下文 & 动作、轨迹 & 次优专家、漏记信息 & 固定日志或另行授权的专家查询 & 行为条件分布；附加理性假设下的奖励约束 \\
    评分 & 单个被评价行为 & 步骤、片段、轨迹 & 尺度漂移、反应延迟 & 可只提供历史评分 & 该尺度下的评价函数 \\
    偏好 & 同条件候选行为 & 一对片段或轨迹 & 次序效应、偏好异质性 & 固定比较或主动选对 & 比较概率、相对效用 \\
    验证 & 产物或状态关系 & 步骤、结果 & 规则遗漏、执行误差、误判 & 可离线复查或在线调用 & 指定条件的通过情况 \\
    \bottomrule
  \end{tabularx}
  \caption{四种反馈形式及其可用信息。形式说明记录表达什么，查询权限说明还能补充什么。人类、AI、规则和执行器描述反馈来源，同一来源可以产生多种形式。}
  \label{tab:rl-reward-forms}
\end{table*}

例如，人和模型都可以提供示范、评分或成对偏好。身份核对规则可以返回通过标记，实际执行开箱操作的检查器也可以返回破损结果。这里“人类反馈”回答谁在判断，“偏好反馈”回答判断怎样记录，两者可以同时描述同一条记录，不是互斥类别。

% Outline: 5.1.2
\subsection{奖励归因受限}
\label{sec:rl-reward-credit-limit}

假设日志已准确记录包裹完好且交付成功，我们仍不知道哪些先前决策值得强化：核对身份、平稳搬运和容易的路线都可能与成功同时出现。\term{奖励归因}（reward credit assignment）要把长期结果联系到先前决策，为局部更新提供贡献判断。第~\ref{chap:rl-learning}章已经用回报、TD误差和优势估计传播结果；这里要检查这些工具接收的信号包含什么信息，以及哪些贡献仍未被观测。

\begin{table*}[htbp]
  \centering\small
  \begin{tabularx}{\linewidth}{@{}>{\raggedright\arraybackslash}p{27mm}*{3}{>{\raggedright\arraybackslash}X}@{}}
    \toprule
    现象 & 缺口 & 主要影响 & 不能直接推出什么 \\
    \midrule
    非零奖励很少 & 多数尝试没有区分信号 & 难以探索到成功轨迹 & 一旦成功就能正确归因 \\
    反馈延迟长 & 行动与反馈相隔许多步 & 贡献需跨步传播；随机轨迹下估计可能更难 & 回报方差必然增大或反馈语义错误 \\
    只有轨迹总分 & 局部贡献未被观测 & 无关动作也可能分享相同结果 & 每一步都有独立监督 \\
    反馈噪声大 & 相同行为的判断不稳定 & 估计方差和误排序增加 & 噪声一定来自时间延迟 \\
    \bottomrule
  \end{tabularx}
  \caption{影响奖励归因的四种反馈性质。非零奖励少、反馈来得晚、只评价整条轨迹及判断有噪声，可以在同一任务中同时出现，但各自造成的缺口不同。}
  \label{tab:rl-reward-credit-limits}
\end{table*}

反馈到达得晚，不等于回报方差必然增大。确定性终局得分无论晚到多少步，方差仍为零。若固定策略和随机得分$X$的分布，只把奖励发生时刻从0移到$h$，在固定$\gamma<1$下，折扣回报的方差为$\gamma^{2h}\operatorname{Var}(X)$，反而不大于原方差。较长交互中，随机状态、动作及评分噪声可能使局部贡献更难估计；判断方差怎样变化时，必须固定所比较的估计量及其他条件。若只是记录晚到、奖励所属时刻不变，离线计算的回报也不会因此改变。

按照表~\ref{tab:rl-reward-credit-limits}，我们可以先区分尚未找到成功行为与已经找到成功行为却不清楚贡献这两种情况。机器人始终找不到收件人时，增加探索或提供成功示范可能补上行为线索。若日志中已有许多成功配送，但身份核对总与无关的扫码动画一起出现，就需要辨认哪一个动作影响结果。重复收集相同组合未必能区分二者；给每次动画播放加分虽然增加了反馈频率，也没有提供更可信的贡献信息。

% Outline: 5.1.3
\subsection{奖励与行为目标不匹配}
\label{sec:rl-reward-mismatch-limit}

若计分规则是每移动一米得一分，机器人可以通过绕圈提高回报，却始终不完成配送。优化器在这种行为上可能正按目标工作，问题出在数值标准没有表达交付要求。要判断这一偏差，我们需要把希望发生的行为与训练实际优化的数值分开。

\begin{definition}[奖励、回报、效用与行为要求]
  \label{def:rl-reward-objective}
  单步奖励$r_t$是交互协议返回的标量，累计回报$G_0=\sum_t\gamma^t r_t$按约定的时间权重汇总这些标量。\term{效用}（utility）表示决策者如何评价结果或结果分布，可记为$u(\tau)$或策略评价泛函$\mathcal U(\pi)$，不预设它等于$J(\pi)$。\term{行为要求}（behavioral requirement）规定希望发生什么以及在什么条件下发生，例如交给正确收件人、保持包裹完好并满足能耗要求。\term{代理奖励}（proxy reward）是训练所用的数值编码；它是否表达了这些要求，须用独立证据检验。
\end{definition}

据此再看三种限制：里程计分即使即时且准确，也可能遗漏任务目标；“正确收件人且完好”的终局标记可以表达明确的成功条件，但早期动作要经过多步才得到反馈。有效路线示范能够帮助行动，却未必确定各项奖励权重。同一任务还可以同时具备三种限制，例如示范不记录破损、评分延迟到终局，训练又额外奖励里程。

选择学习办法要先定位缺口。缺少可靠信号时需要补充目标信息，缺少局部贡献时需要检查结果怎样进入更新，标准遗漏要求时则需要修订目标。训练损失下降只能说明模型更符合所选训练准则，不能同时证明反馈覆盖充分、归因正确和编码完整。交付身份、完好情况及能耗的独立检查应在优化前约定；若事后随行为改变标准，就失去了判定偏差的共同依据。

% Outline: 5.2
\section{奖励观测受限下的强化学习}
\label{sec:rl-reward-alternative}

假设配送日志拿不到可靠的逐步奖励，却保留了专家路线、两条路线的比较或签收检查。我们可以用动作示范直接拟合策略，用偏好约束相对效用，也可以用验证结果形成策略更新信号，不必都先恢复唯一奖励。下面按记录实际表达的内容，判断能从中推断什么，以及推断还依赖哪些假设。

% Outline: 5.2.1
\subsection{替代反馈与可推断对象}
\label{sec:rl-reward-inference}

在同一个配送任务中，示范可以记录专家选择电梯，偏好可以记录收件人更喜欢平稳但较慢的配送，身份验证则只检查是否交给正确的人。选择电梯没有单独说明“慢一些更好”，身份通过也没有证明包裹无破损。若把这些记录统一改成“正例”，我们就无法再分清它们支持哪种结论、排除了哪种目标解释。

为了说明这些结论的对象，我们回用前文的记号。策略$\pi(a\mid s)$给出某状态上的动作分布，占用度量$d^\pi(s,a)$汇总长期访问状态和动作的概率；价值$V^\pi,Q^\pi$预测给定奖励和策略下的后续回报。奖励$R(s,a)$规定单步期望计分，效用表达评价者怎样看待结果。这些对象相互联系，但一条记录对其中一个对象提供信息，并不自动唯一确定其他对象。表~\ref{tab:rl-reward-inference}列出各类记录可以支持的学习对象。

\begin{table*}[htbp]
  \centering\small
  \begin{tabularx}{\linewidth}{@{}>{\raggedright\arraybackslash}p{22mm}*{4}{>{\raggedright\arraybackslash}X}@{}}
    \toprule
    反馈 & 候选推断对象 & 额外假设 & 训练交互 & 最终输出 \\
    \midrule
    状态动作示范 & 策略 & 输入足以解释动作，覆盖可用 & 行为克隆不需要 & 动作分布 \\
    专家与策略轨迹 & 占用、软价值 & 共同环境，轨迹采样口径一致 & 占用匹配通常需要 & 策略或价值函数 \\
    专家轨迹 & 奖励、特征期望 & 专家理性、奖励函数类、动态模型 & 需规划或采样权限 & 奖励候选及策略 \\
    标量评分 & 评价函数 & 尺度、对象与时间归属明确 & 拟合本身不需要 & 评分预测器 \\
    成对偏好 & 效用差、策略 & 比较模型与共同条件 & 拟合可离线；在线RL另需新采样 & 奖励模型或偏好策略 \\
    过程、结果验证 & 条件通过概率 & 规则有效、检查前提满足 & 在线更新需可调用验证 & 验证信号与策略 \\
    \bottomrule
  \end{tabularx}
  \caption{不同记录可支持的推断、假设和输出。同一反馈可以用于多个学习对象，同一对象也可以结合多种反馈。“不需要交互”只说明拟合阶段的访问要求，不表示固定数据已覆盖未来行为。}
  \label{tab:rl-reward-inference}
\end{table*}

一个单步例子即可看出不唯一性。任务只有左、右两个动作，专家总选左；奖励$(R(\text{左}),R(\text{右}))=(1,0)$和$(10,9)$都能使左最优。若记录只保留甲、乙、丙的确定排序，效用$(3,2,1)$和$(100,2,0)$也都给出甲优于乙、乙优于丙，但它们的效用差不同，未必在同一个随机偏好模型中产生相同比较概率。排序提供的信息比完整比较概率少；即使概率已知，共同加法常数等歧义通常仍然存在。

新增查询能否解除歧义，取决于它提出了什么新问题。若允许主动查询，我们可以比较此前未比较的路线，或询问专家在当前策略偏离后怎样恢复。重复旧记录可以降低采样误差，却不会补出从未涉及的要求。因而这里沿用第~\ref{chap:rl-information}章的判断顺序：先说明要识别的是动作、排序还是奖励，再检查数据是否足以支持这个结论。

% Outline: 5.2.2
\subsection{从示范直接学习策略}
\label{sec:rl-reward-bc-dagger}

专家已经给出可用动作时，我们可以先让策略在相似输入上复现这些动作。前章用于拟合日志动作的\term{行为克隆}（behavioral cloning, BC），在这里以专家状态或记录历史为输入，以专家动作作监督标签。要衡量离散动作策略是否复现了示范，可以惩罚它给专家动作分配过小的概率。设任务有$t=0,\ldots,H-1$这$H$步，专家在各时刻的状态分布为$\nu_t^{\mathrm E}$；对这些状态及专家动作取期望，再在时间上平均，得到总体负对数似然
\begin{equation}
  L_{\mathrm{BC}}(\bm{\theta})
  =-\frac1H\sum_{t=0}^{H-1}
  \mathbb E_{\substack{s\sim\nu_t^{\mathrm E}\\
              a\sim\pi_{\mathrm E}(\cdot\mid s)}}
       \log\pi_{\bm{\theta}}(a\mid s).
  \label{eq:rl-reward-bc}
\end{equation}
最小化该损失，就是在专家访问分布上最大化示范动作的平均对数概率，并不要求每条记录的概率都同时增加。训练时用示范数据$\mathcal D_{\mathrm E}$的平均负对数似然估计这个期望，再做
$\bm{\theta}^{(k+1)}=\bm{\theta}^{(k)}
-\eta\nabla_{\bm{\theta}}\widehat L_{\mathrm{BC}}$，
完成一次参数拟合，其中$\eta$为步长。数据由专家产生，不是当前策略执行得到的；每次遍历都要为记录计算策略概率和梯度，但不需要奖励标签或重新运行环境。

连续动作还可以用平方误差$\mathbb E\lVert f_{\bm{\theta}}(s)-a\rVert_2^2$拟合动作数值；在无限容量下，最小化这个误差得到条件均值。然而，均值未必是可用动作。假设专家绕障时等概率选择$-1$和$1$，分别表示左转和右转，均值$0$可能使机器人直撞障碍。要保留两种有效选择，就需表示多模态动作分布，而不只输出平均动作。另一个原因是日志漏记了任务指令，使相同观测对应不同专家动作；单靠增加网络复杂度不能恢复未记录的指令。

动作在专家数据上拟合准确，还不能保证执行时同样准确。专家很少进入错误走廊，克隆策略一旦误入，就要在示范没有覆盖的状态继续决策。为了看清早期错误如何影响后续成本，下面在共同环境中比较专家与学习策略，并用首次动作分歧的概率给出保守上界。

\begin{theorem}[行为克隆的首次分歧界]
  \label{thm:rl-reward-bc}
  考虑共同初始分布、共同转移核的有限时域任务，每步成本$c_t(s,a)\in[0,1]$，总成本记为
  $C_H(\pi)=\mathbb E_\pi\sum_{t=0}^{H-1}c_t(s_t,a_t)$。
  专家$\pi_{\mathrm E}$为确定性策略。若学习策略在每个专家状态分布上均满足
  \[
    \mathbb E_{s\sim\nu_t^{\mathrm E}}
    \bigl[1-\pi(\pi_{\mathrm E}(s)\mid s)\bigr]\leq\varepsilon,
  \]
  则
  \[
    C_H(\pi)-C_H(\pi_{\mathrm E})
    \leq \frac{H(H+1)}2\varepsilon.
  \]
  非平稳专家可把时间加入状态，结论不变。
\end{theorem}

\begin{proof}
  将两条轨迹放在同一概率空间中，使它们从相同状态出发，并在尚未出现动作分歧时共用环境转移的随机性。另在完整专家轨迹的每个$s_t^{\mathrm E}$上，抽取服从$\pi(\cdot\mid s_t^{\mathrm E})$的候选动作。令$B_t$表示该候选动作不同于专家动作；定理假设给出$\mathbb P(B_t)\leq\varepsilon$。

  首次分歧以前，两条轨迹的状态、动作和成本相同，所以截至$t$已发生分歧，只可能是在某个$j\leq t$出现了$B_j$。这个事件包含于$\bigcup_{j=0}^{t}B_j$，由并集界，其概率不超过$(t+1)\varepsilon$。分歧后两条轨迹可以任意发展，但单步成本均在$[0,1]$内，额外成本至多为1；未分歧时额外成本为0。对这一指示变量上界取期望，再按时间求和，
  \[
    C_H(\pi)-C_H(\pi_{\mathrm E})
    \leq\sum_{t=0}^{H-1}(t+1)\varepsilon
    =\frac{H(H+1)}2\varepsilon.
  \]
  并集界使用的是完整专家轨迹上的事件概率，并未假定“此前尚未分歧”条件下的错误率也等于$\varepsilon$，因此没有额外引入条件分布假设。
\end{proof}

由于每步成本至多为1，总成本差本身还不超过$H$，可以将界的右端取为$H$与上述二次项中的较小值。二次项来自一次早期分歧可能影响余下多步的估计，不是所有任务都会二次恶化的断言；若机器人能迅速恢复，实际额外成本可以小得多。

上述缺口发生在学习者实际访问的状态上，因此可以让专家对这些状态补动作标签。\term{数据集聚合}（Dataset Aggregation, DAgger）反复执行当前策略、查询专家，并把新标签并入训练集重新拟合。它改变训练覆盖，使偏离后的状态也进入下一轮，而不是只增加专家原有轨迹。图~\ref{fig:rl-reward-dagger}分别画出执行、查询和聚合；Ross等给出了这一归约及其在线学习分析\scite{ross2011}

\begin{algorithm}[一轮DAgger更新]
  \label{alg:rl-reward-dagger}
  \begin{algorithmic}[1]
    \Require 已聚合数据$\mathcal D^{(k)}$，策略$\pi^{(k)}$，专家查询接口，轨迹预算$m$
    \Ensure 新数据$\mathcal D^{(k+1)}$与候选策略$\pi^{(k+1)}$
    \State 在获准环境中执行$m$条$\pi^{(k)}$轨迹，记录实际访问的状态
    \State 对这些状态查询$a_{\mathrm E}=\pi_{\mathrm E}(s)$，保留失败及偏离后的状态
    \State $\mathcal D^{(k+1)}\gets\mathcal D^{(k)}\cup\{(s,a_{\mathrm E})\}$
    \State 在聚合数据上最小化模仿损失，得到$\pi^{(k+1)}$
    \State 在独立轨迹上评价任务表现，并记录查询数与执行风险
  \end{algorithmic}
\end{algorithm}

\begin{figure}[htbp]
  \centering
  % Paths are schematic, not trajectories from the stochastic grid.
\begin{tikzpicture}[
  x=1mm,y=1mm,font=\figurefont\fontsize{8}{10}\selectfont,
  state/.style={circle,draw=BookBlue,fill=white,minimum size=8mm,
    inner sep=1mm,line width=.8pt},
  flow/.style={-{Stealth[length=2mm]},draw=BookBlue,line width=.9pt},
  query/.style={-{Stealth[length=2mm]},draw=BookAmber,dashed,line width=.8pt},
  box/.style={draw=BookRule,fill=BookPaper,align=center,inner sep=2mm}
]
  \node[state] (s0) at (7,0) {$s_0$};
  \node[state] (e1) at (38,15) {$s_1^{\mathrm E}$};
  \node[state] (e2) at (72,15) {$s_2^{\mathrm E}$};
  \node[state] (p1) at (38,-7) {$s_{\text{偏}}$};
  \node[state] (p2) at (72,-7) {$s_{\text{新}}$};
  \draw[flow] (s0) -- node[above,sloped] {专家执行} (e1);
  \draw[flow] (e1) -- (e2);
  \draw[flow,draw=BookTeal] (s0) -- node[below,sloped] {当前策略执行} (p1);
  \draw[flow,draw=BookTeal] (p1) -- (p2);
  \node[box] (expert) at (39,-29) {专家接口\\返回动作标签};
  \node[box] (data) at (85,-29) {聚合数据\\$\mathcal D^{(k+1)}$};
  \draw[query] (p1) -- node[left] {查询} (expert);
  \draw[query] (p2.south) -- (72,-18) -| (expert.east);
  \draw[query] (expert) -- node[above] {聚合} (data);
  \node[box] (fit) at (85,-50) {重新拟合$\pi^{(k+1)}$};
  \draw[flow,draw=BookMuted] (data) -- node[right] {监督学习} (fit);
  \node[align=left,anchor=west] at (0,-49) {新状态进入训练\\不只重放专家旧轨迹};
\end{tikzpicture}

  \caption{DAgger补充当前策略偏离后的状态。两条轨迹从同一起点出发，分歧后对当前策略访问的新状态查询专家；新标签随聚合数据进入下一轮。查询只取得动作标签，不表示专家已在环境中执行该动作。}
  \label{fig:rl-reward-dagger}
\end{figure}

例如，当前策略误入$s_{\text{偏}}$，专家标注“返回”。在二动作策略中，用logit $z$表示返回的对数几率，返回概率为$\sigma(z)$，其中$\sigma(z)=(1+\exp(-z))^{-1}$将实数转换为概率。若当前$z=0$，该标签的负对数似然梯度是$\sigma(0)-1=-1/2$；以$0.2$为步长下降，得到$z=0.1$，返回概率由$0.5$升至约$0.525$。这次更新利用了原专家示范未覆盖的错误状态。取得这些标签也有成本：每轮执行$m$条$H$步轨迹需要$mH$个环境步、至多$mH$次查询，并要重新拟合不断增长的数据集。

能否持续改善自身访问分布上的模仿损失，可以用在线学习的遗憾来描述：每轮损失在该轮策略访问分布上计算，比较基准则是预先指定策略类$\Pi$中，同一策略在全部轮次上的平均损失。下面先假定这些总体损失可用；平均遗憾小意味着迭代策略的平均损失接近类内最小值，因而至少有一轮不超过该平均上界。

\begin{theorem}[理想化DAgger的无遗憾归约]
  \label{thm:rl-reward-dagger}
  令$\bar\nu^{(k)}=H^{-1}\sum_t\nu_t^{\pi^{(k)}}$为第$k$轮当前策略自身的平均状态分布，定义
  $\ell_k(\pi)=\mathbb E_{s\sim\bar\nu^{(k)}}\ell(\pi;s)$。
  若$K$轮在线拟合满足
  \[
    \frac1K\sum_{k=1}^K\ell_k(\pi^{(k)})
    -\min_{\pi\in\Pi}\frac1K\sum_{k=1}^K\ell_k(\pi)
    \leq\varepsilon_{\mathrm{reg}},
  \]
  则至少一轮$k$满足
  \[
    \ell_k(\pi^{(k)})
    \leq\min_{\pi\in\Pi}\frac1K\sum_{j=1}^K\ell_j(\pi)
        +\varepsilon_{\mathrm{reg}}.
  \]
\end{theorem}

将遗憾不等式中的比较项移到右侧，就得到各轮自身分布损失的平均上界。再在实际生成的$K$个策略中取损失最小的一轮，利用
\[
  \min_k\ell_k(\pi^{(k)})
  \leq\frac1K\sum_k\ell_k(\pi^{(k)}).
\]
即可得到定理的结论。比较策略在所有轮次的分布上评价，分布本身不必相同。这个结果首先约束模仿损失；使用有限样本还要加入估计误差，混合执行专家时还要控制混合分布与自身分布的差异。若进一步推成任务成本界，需要模仿损失能控制动作错误，并且一次错误后专家恢复的额外成本有统一上界。平均遗憾条件本身不足以保证任意实现都具有线性时域成本界。

仅有动作标签时，DAgger不能直接区分两种错误动作的后续代价。基于搜索的结构化预测（search-based structured prediction, SEARN）将序贯预测转为代价敏感学习，通常通过后续执行估计各动作的代价；其中SEARN是方法专名，不是该英文名称逐词首字母的缩写。聚合价值模仿（Aggregate Values to Imitate, AggreVaTe）则聚合当前访问状态上的专家后续成本，使“动作不同但成本接近”也能进入比较。它们比动作查询需要更丰富的信息和更高的获取成本。专家还可能无法评价陌生状态，采集时的策略执行也可能造成风险；读取固定示范不意味着已经获准执行这些查询\scite{rl5-daume2009searn,rl5-ross2014aggrevate}

% Outline: 5.2.3
\subsection{从示范推断占用与奖励}
\label{sec:rl-reward-occupancy-irl}

配送策略可能在常见走廊上选择与专家相同的动作，却在装载区停留更久；只拟合专家已见状态上的动作，尚未比较这种长期访问差异。我们回用归一化折扣占用度量，将不同时刻的状态动作分布按折扣汇总。这里改用$0\leq\gamma<1$的无限折扣设定，回合结束后以原始奖励为零的吸收状态延拓，并将终止后的占用计入总质量：
\begin{equation}
  d^\pi(s,a)=(1-\gamma)\sum_{t=0}^{\infty}
       \gamma^t\mathbb P_\pi(s_t=s,a_t=a).
  \label{eq:rl-reward-occupancy}
\end{equation}
该分布同时记录策略怎样选动作以及环境允许到达哪些状态。在共同环境中，完整占用匹配不仅要求已见状态上的动作相符，还约束这些状态出现的频率。对于平稳马尔可夫策略，在$d^\pi(s)>0$处可以由$d^\pi(s,a)/d^\pi(s)$恢复动作分布。比较前必须将专家和策略样本按同一折扣及吸收延拓方式加权；若在不同长度回合中均匀抽取非终止转移，通常得到的是另一种采样分布，不是这里的$d^\pi$。

\term{生成对抗模仿学习}（generative adversarial imitation learning, GAIL）用判别器识别状态动作来自专家还是当前策略，再更新策略，使两种来源更难区分。令$D_{\bm{\omega}}(s,a)$为判别器给出的专家来源概率。判别器在允许的概率函数中最大化正确来源的对数似然，策略在共同环境可实现的策略中最小化这一可区分程度，并以非负权重$\lambda$鼓励动作熵。使用归一化占用，可写为
\begin{equation}
  \min_\pi\max_D
  \left\{\mathbb E_{d^{\mathrm E}}\log D
    +\mathbb E_{d^\pi}\log(1-D)
    -\lambda\mathcal H_d(\pi)\right\},
  \label{eq:rl-reward-gail}
\end{equation}
其中$\mathcal H_d(\pi)=\mathbb E_{d^\pi}[-\log\pi(a\mid s)]$在策略访问分布上计算动作熵。若暂不限制判别器函数类，固定两份占用，在有总质量的位置逐点最大化判别项，得到
$D^*=d^{\mathrm E}/(d^{\mathrm E}+d^\pi)$；
代回判别项，得到$-\log4+2\operatorname{JS}(d^{\mathrm E}\Vert d^\pi)$。这里$\operatorname{JS}$是Jensen--Shannon散度，即分别计算两个分布相对于它们平均分布的KL，再取平均；值越小表示访问分布越难区分。这个推导说明判别器提供的是占用差异的度量，不是唯一真实奖励的识别结果\scite{rl5-ho2016gail}

实际训练先用当前策略采集轨迹，将专家与策略样本分别标为1和0，通过交叉熵拟合判别器。随后冻结判别器，以
$\widetilde r_t=-\log(1-D_{\bm{\omega}}(s_t,a_t))$
构造策略优化的奖励：最大化这个负对数项，等价于降低式~\eqref{eq:rl-reward-gail}中的策略判别项。若两个策略样本分别得到$D=0.8$和$D=0.2$，奖励约为$1.609$和$0.223$，第一项更鼓励模仿专家，却没有证明它满足示范未记录的配送要求。若优化器使用未归一化回报，还要按前文约定换算为
$\widetilde J(\pi)=(1-\gamma)^{-1}\mathbb E_{d^\pi}\widetilde R$，
其中$\widetilde J$只表示判别奖励的回报。将整个目标乘同一正常数不改变最优策略，但会改变梯度尺度；若只换算奖励项，还必须相应调整正则化权重，才能保留两项之间的取舍。

为了比较当前策略的占用，通常需要继续与环境交互，专家一侧则可以使用固定轨迹。新轨迹采集和判别器、策略的交替优化构成主要成本，判别器过拟合、两份占用支持几乎不交叠或策略更新滞后，都可能妨碍匹配。即使占用相同，也未必满足全部轨迹要求；例如要检查步骤次序，可能需在状态中加入此前步骤的信息。

占用判别并不是从示范学习的唯一入口。若希望避免独立奖励网络与策略网络的对抗更新，可以直接学习软动作价值，即计入动作熵后评价后续回报的$Q$。\term{逆软Q学习}（Inverse soft-Q Learning, IQ-Learn）用软Bellman关系将奖励与策略都写成同一个$Q$的函数。在有限离散动作、熵温度取1时，先对合法动作的指数评分求和再取对数，得到软状态价值，再据此归一化动作概率，并由Bellman残差表示隐含奖励：
\begin{equation}
  \begin{aligned}
    V_Q(s)&=\log\sum_a\exp Q(s,a),\\
    \pi_Q(a\mid s)&=\exp\{Q(s,a)-V_Q(s)\},\\
    R_Q(s,a)&=Q(s,a)-\gamma\mathbb E_{s'\mid s,a}V_Q(s').
  \end{aligned}
  \label{eq:rl-reward-iq}
\end{equation}
式~\eqref{eq:rl-reward-iq}让一个$Q$同时导出动作分布和隐含奖励，但还没有规定怎样从专家数据选择$Q$。为使专家行为在隐含奖励下取得较高评价，同时限制奖励幅度，IQ-Learn采用专家奖励项与正则项组成的目标。这里选原论文讨论的$\chi^2$正则形式，保持上述温度1及归一化占用。令$\rho_0$为共同初始状态分布，$\kappa>0$为尺度参数，$c=1/(4\kappa)$为奖励平方项的系数；在选定的有界$Q$函数类内最大化
\[
  \mathcal J_{\mathrm{IQ}}(Q)
  =\mathbb E_{d^{\mathrm E}}
       \left[R_Q(s,a)-cR_Q(s,a)^2\right]
     -(1-\gamma)\mathbb E_{s_0\sim\rho_0}V_Q(s_0).
\]
第一项提高专家访问位置上的隐含奖励，平方项抑制这些位置的奖励幅度，不是对策略样本奖励平方作正则。最后一项扣除软价值的初态评价，抵消整体抬高$Q$带来的线性奖励增量；整体平移仍可能改变平方正则项。由软Bellman关系，
$(1-\gamma)\mathbb E_{\rho_0}V_Q
=\mathbb E_{d^{\pi_Q}}R_Q+\mathcal H_d(\pi_Q)$；
因此目标比较专家与$\pi_Q$的隐含奖励，同时扣除上述策略熵项；优化对象仍是决定奖励及策略的$Q$，不能将它当作固定奖励下的策略目标。这一平方正则来自$\chi^2(d^{\pi_Q}\Vert d^{\mathrm E})$占用差异的对偶形式。在离散状态动作上，该散度是
$\sum_{s,a}[d^{\pi_Q}(s,a)-d^{\mathrm E}(s,a)]^2/d^{\mathrm E}(s,a)$，
比较的是策略占用相对专家占用，不是“恢复真实奖励”的误差。散度对应完整的对偶函数范围，受限$Q$类下采用的是相应受限目标。若策略占用超出专家支持，这种散度可以无穷大；有限数据与受限函数类下的目标不能自动排除该问题\scite{rl5-garg2021iq}

为了给出一次可执行的更新，先看只用专家转移的离线形式。同一折扣与吸收延拓下，占用流量恒等式给出
$\mathbb E_{d^{\mathrm E},P}[V_Q(s)-\gamma V_Q(s')]
=(1-\gamma)\mathbb E_{\rho_0}V_Q(s_0)$，
其中$P$为真实转移核。设批次$\{(s_i,a_i,s'_i)\}_{i=1}^n$按专家归一化占用及转移采样，记单样本残差
$\widehat\delta_i=Q(s_i,a_i)-\gamma V_Q(s'_i)$。
用这些样本构造原论文离线形式的经验损失，并下降更新：
\[
  \widehat L_{\mathrm{IQ}}(\bm\theta)
  =\frac1n\sum_{i=1}^n
      \left[V_{Q_{\bm\theta}}(s_i)-Q_{\bm\theta}(s_i,a_i)
                    +c\widehat\delta_i^2\right],
  \qquad
  \bm\theta^{(k+1)}
  =\bm\theta^{(k)}-\eta\nabla_{\bm\theta}\widehat L_{\mathrm{IQ}}.
\]
这里$\eta>0$为步长，$\widehat\delta_i$也由$Q_{\bm\theta}$计算；所示梯度针对这一经验损失。原始轨迹可按几何折扣抽样或用等价权重，不能把任意回合的非终止转移均匀平均冒充该占用。确定性转移下$\widehat\delta_i=R_Q(s_i,a_i)$；随机转移下，条件均值仍为$R_Q$，但
$\mathbb E[\widehat\delta_i^2\mid s_i,a_i]
=R_Q(s_i,a_i)^2+\gamma^2\operatorname{Var}[V_Q(s'_i)\mid s_i,a_i]$。
故单后继平方项是经验代理，不是总体奖励平方项的无偏估计。吸收状态的唯一虚拟动作没有动作熵，仍须按共同协议计入占用；环境原始奖励在那里为零，不表示学出的$R_Q$必须强置零。

在线版本还可以采集策略转移，用它们估计软价值项；离线版本不需要刷新策略轨迹，却仍受固定数据支持限制。连续动作通常需要参数化行动者近似软策略，而不再直接枚举动作求和。减少独立网络的对抗循环，也不意味着$R_Q$已具有唯一的奖励语义。

若我们要解释专家为何选择这些动作，而不只复现行为，就需要约束奖励本身。\term{逆强化学习}（inverse reinforcement learning, IRL）寻找能够支持专家选择的奖励。Ng与Russell把专家最优性写成奖励约束：已知模型和专家策略时，对每个备选动作要求$Q_R^{\pi_{\mathrm E}}(s,\pi_{\mathrm E}(s))\geq Q_R^{\pi_{\mathrm E}}(s,a)$。这表示遵循专家的后续策略时，专家动作不比备选动作差。但零奖励也满足所有不等式，使全部动作并列最优；要缩小候选范围，还需规定奖励函数类、幅度约束或动作价值间隔。专家若只是近似最优，容许误差也必须明示\scite{ng2000irl}

无法恢复唯一奖励时，还可以问策略在某一奖励类中是否具有接近专家的回报。学徒学习（apprenticeship learning）采用这一较弱目标，用示范约束长期行为特征。假设奖励为$R_{\bm w}(s,a)=\bm w^\top\bm\phi(s,a)$，其中$\bm\phi$为选定的状态动作特征，$\bm w$为权重。定义归一化特征期望$\bm m_\pi=\mathbb E_{d^\pi}\bm\phi(s,a)$后，利用占用与未归一化回报的关系可得
\begin{equation}
  J_{\bm w}(\pi)-J_{\bm w}(\pi_{\mathrm E})
  =\frac{\bm w^\top(\bm m_\pi-\bm m_{\mathrm E})}{1-\gamma}.
  \label{eq:rl-reward-feature}
\end{equation}
若权重满足$\lVert\bm w\rVert_2\leq1$，且两份特征期望的欧氏距离不超过$\delta$，Cauchy--Schwarz不等式给出绝对回报差至多$\delta/(1-\gamma)$。匹配特征因此可以同时控制这类线性奖励的回报差，却无法保证特征未编码的破损要求。计算候选策略的特征期望还需要模型规划或环境采样权限\scite{rl5-abbeel2004apprentice}

专家也可能在多条可行路线间变化选择，单纯要求一个动作最优不足以描述示范频率。最大熵逆强化学习用轨迹概率模型解释这些变化：较高累计奖励的轨迹取得较大概率，同时保留选择的不确定性。在固定起点、有限时域及确定性动力学的简化设定下，以环境允许的轨迹为候选集合，取
\[
  p_{\bm w}(\tau)=\frac{\exp\{\bm w^\top\bm f(\tau)\}}{Z(\bm w)},
  \qquad \bm f(\tau)=\sum_t\gamma^t\bm\phi(s_t,a_t).
\]
$\bm f(\tau)$是整条轨迹的折扣特征和，$Z(\bm w)$对所有可行轨迹的指数权重求和，使概率归一化。最大化示范平均对数似然，其梯度为
$\mathbb E_{\mathcal D_{\mathrm E}}\bm f-\mathbb E_{p_{\bm w}}\bm f$，
因此更新增加示范中出现更多、模型中出现更少的特征权重。配分函数和模型期望的计算需要动态规划或采样近似。在随机动力学下，直接给完整未来轨迹加权可能连环境噪声概率也一并改变，而智能体只能选择动作，不能选择噪声。最大因果熵表述据此保留真实转移，只优化依据已有信息可执行的动作分布\scite{rl5-ziebart2008maxent}

\begin{table*}[htbp]
  \centering\small
  \begin{tabularx}{\linewidth}{@{}>{\raggedright\arraybackslash}p{25mm}*{4}{>{\raggedright\arraybackslash}X}@{}}
    \toprule
    路线 & 观测输入 & 学习对象 & 输出 & 所需权限 \\
    \midrule
    直接模仿 & 专家状态动作 & 条件动作分布 & 策略 & BC可只用日志；DAgger另需执行与查询 \\
    占用匹配 & 专家及策略轨迹 & 占用差异、判别器 & 策略 & 通常持续采样环境 \\
    软价值学习 & 专家及可用转移 & 软$Q$及隐含奖励 & 价值、策略 & 离线或在线版本分开 \\
    特征期望匹配 & 示范特征与模型行为 & 特征期望差 & 一类奖励下可用的策略 & 规划或采样候选策略 \\
    奖励推断 & 专家行为与理性模型 & 奖励参数 & 奖励候选及对应策略 & 动态模型或可估计的转移信息 \\
    \bottomrule
  \end{tabularx}
  \caption{示范学习路线的输入、输出及权限。各方法学习的对象不同；即使算法输出奖励候选，也不表示数据已唯一确定该奖励。}
  \label{tab:rl-reward-imitation}
\end{table*}

表~\ref{tab:rl-reward-imitation}应按所求结论来读，而不只按方法名称选择。专家选电梯既可能看重每步能耗低，也可能看重包裹保持水平；两种奖励都支持当前动作，却可能在电梯停运后选择不同替代路线。因此，恢复原环境的可用策略，比恢复可跨环境使用的奖励要求更弱。第~\ref{sec:rl-reward-ambiguity}节将用未比较路径具体检验这一区别。

% Outline: 5.2.4
\subsection{从评分与偏好学习奖励或策略}
\label{sec:rl-reward-preference}

评价者不便给出动作示范时，可能仍能为机器人刚才的行为打分。通过评价强化手动训练智能体（Training an Agent Manually via Evaluative Reinforcement, TAMER）用人工标量反馈拟合评价函数，再按预测评分选择动作。这里拟合的是人的评价，不自动等于环境长期回报；人的反应还有延迟，必须另行判断评分对应哪次动作，不能按按键到达时刻直接归属反馈\scite{rl5-knox2009tamer}

若难以统一“这一段值几分”的尺度，可以让评价者在同一包裹和路线条件下比较两段配送。为了从带噪声的选择记录学习评价函数，我们需要把效用差与选择概率联系起来。令$x,y$为两个片段，$u_{\bm{\omega}}(x)$为片段效用，$\sigma(z)=(1+\exp(-z))^{-1}$。\term{Bradley--Terry模型}假定效用差越大，$x$胜出的概率越高，并用下式给出具体关系：
\begin{equation}
  \mathbb P_{\bm{\omega}}(x\succ y)
    =\sigma\!\left(u_{\bm{\omega}}(x)-u_{\bm{\omega}}(y)\right).
  \label{eq:rl-reward-bt}
\end{equation}
这种关系是关于比较噪声的可检验假设，不是所有偏好都必然满足的定义。原模型用两个正强度的相对大小表示成对选择概率，将强度取对数，就得到这里的效用差形式\scite{rl5-bradley1952paired}

拟合时，应让实际被选对象取得较大的预测选择概率。用$b=1$记录选$x$，用$b=0$记录选$y$，并记$\Delta_j=u_{\bm{\omega}}(x_j)-u_{\bm{\omega}}(y_j)$。对$n$个条件独立的比较记录，最小化平均负对数似然
\begin{equation}
  L_{\mathrm{pref}}(\bm{\omega})
  =-\frac1n\sum_{j=1}^n
    \left[b_j\log\sigma(\Delta_j)
      +(1-b_j)\log\sigma(-\Delta_j)\right],
  \label{eq:rl-reward-bt-loss}
\end{equation}
就能以比较标签更新效用参数。损失对$\Delta_j$的导数为$(\sigma(\Delta_j)-b_j)/n$：预测选择概率低于实际标签时，下降更新会增大相应效用差，反之会减小。给两边效用加同一常数不改变似然，所以这个常数不可由比较识别；若同时拟合噪声温度，效用尺度还可能与温度混淆。

要将片段评价用于逐步决策，还需约定效用怎样从局部奖励组成。一种参数化为
$u_{\bm{\omega}}(x)=\sum_{t=0}^{L_x-1}\gamma^t
\widehat R_{\bm{\omega}}(s_t,a_t)$，
其中$L_x$为片段长度，片段起点重新编号为0。这个局部和是附加建模假设，不是成对偏好本身提供的信息。片段长度不同时，比较总量与比较平均质量采用不同标准，改用平均分就改变了效用模型。标注者差异可以用分组或条件变量表示，呈现顺序也应随机化并记录。若稳定偏好形成循环，固定标量效用的严格排序无法表达它；若只是偶尔出现循环标签，则可能是比较噪声。两种情况需分别诊断，不能一律删掉不一致标签。

Deep RL from Human Preferences将片段比较、奖励拟合和策略更新接成循环：奖励模型为策略提供信号，更新后的策略又产生新的待比较片段。相应的语言模型训练常称为\term{基于人类反馈的强化学习}（reinforcement learning from human feedback, RLHF），即将人类评价学成奖励，再用它优化生成策略。经典安排先以监督微调取得初始策略，再用偏好拟合奖励模型，随后优化带参考策略约束的奖励，并独立评价新行为。奖励模型拟合准确与策略行为符合要求，是两个需要分别检验的结论\scite{christiano2017}

在语言模型中，完整回答是被评价对象，词元（token）则是每步生成的文本单位，作为策略动作，二者粒度不同。令$q$为输入，$y=(y_0,\ldots,y_{T-1})$为含结束符的回答，生成前缀构成$s_t=(q,y_{<t})$，动作是$a_t=y_t$。这里采用最大长度内自然结束的有限序列，以未折扣总分评价，相当于$\gamma=1$，不使用无限折扣占用公式；强制截断时的评分及续接规则需要另行规定。奖励模型$\widehat r_{\bm{\omega}}(q,y)$为完整回答输出一个分数，价值模型$V_{\bm{\psi}}(s_t)$预测前缀之后的回报，策略$\pi_{\bm{\theta}}$逐token采样。图~\ref{fig:rl-reward-rlhf}区分这三种计算：一个序列分数可以传播到各步更新，却没有增加逐token独立标签。InstructGPT给出了经典语言模型RLHF流程\scite{ouyang2022instructgpt}

为更新这批token动作，需要将终局评分与偏离参考策略的代价放入同一回报。冻结当前轮旧策略$\mu=\pi^{(k)}$、奖励模型和参考策略$\pi_{\mathrm{ref}}$，在旧策略样本上构造
\begin{equation}
  \widetilde r_t
  =\mathbf1\{t=T-1\}\widehat r_{\bm{\omega}}(q,y)
  -\beta\log\frac{\mu(a_t\mid s_t)}
                      {\pi_{\mathrm{ref}}(a_t\mid s_t)}.
  \label{eq:rl-reward-token-reward}
\end{equation}
其中$\beta>0$控制参考偏移的惩罚。旧策略和参考策略在各允许前缀的共同合法token集合上均须为正，新策略也不能超出该支持。采样对数概率比可以为负；只有在给定前缀并按$\mu$对动作取期望时，它才等于非负的前向KL。按照自然结束时$V_{\bm{\psi}}(s_T)=0$的边界条件，计算
$\delta_t=\widetilde r_t+V_{\bm{\psi}}(s_{t+1})-V_{\bm{\psi}}(s_t)$
及$\widehat A_t=\sum_{j=t}^{T-1}\lambda^{j-t}\delta_j$。
其中$0\leq\lambda\leq1$调节未来TD误差的权重，$\widehat A_t$用来判断当前token相对基线是否值得增加概率。令$\varrho_t(\bm{\theta})=\pi_{\bm{\theta}}(a_t\mid s_t)/\mu(a_t\mid s_t)$，并以$\epsilon>0$指定概率比剪裁范围。PPO式更新在参数化策略类中最大化下面各回答求和项的样本平均：
\begin{equation}
  \sum_{t=0}^{T-1}
  \min\!\left\{\varrho_t\widehat A_t,\,
  \operatorname{clip}(\varrho_t,1-\epsilon,1+\epsilon)\widehat A_t\right\}.
  \label{eq:rl-reward-ppo}
\end{equation}
取两项中的较小值，限制的是样本代理目标对概率比变化的收益，不是为真实策略施加硬概率约束。价值参数另用回报目标回归；本次策略更新将奖励和优势视为固定，下一轮才刷新旧策略与轨迹。参考KL也可以放进显式损失，但需使用配套推导，避免在奖励和损失中重复计罚。序列采样、参考概率计算、奖励打分及价值拟合，都是这一安排的计算成本。

\begin{figure*}[htbp]
  \centering
  % Data and update signals; no gradient through discrete sampling.
\begin{tikzpicture}[
  x=1mm,y=1mm,font=\figurefont\fontsize{8}{10}\selectfont,
  box/.style={draw=BookRule,fill=BookPaper,text width=26mm,
    minimum height=11mm,align=center,inner sep=2mm,line width=.7pt},
  flow/.style={-{Stealth[length=2mm]},draw=BookMuted,line width=.8pt},
  lab/.style={fill=white,inner sep=.7mm,font=\figurefont\fontsize{7}{9}\selectfont}
]
  \node[box,fill=FigureInput!10] (input) at (16,32) {输入$q$};
  \node[box,fill=FigureModel!10] (policy) at (61,32) {旧策略$\mu=\pi^{(k)}$\\逐token采样};
  \node[box] (response) at (107,32) {回答$y$、前缀$s_t$\\保留采样概率$\mu$};
  \node[box,fill=FigureLoss!10] (reward) at (145,32) {奖励模型\\$\widehat r_{\bm\omega}(q,y)$};
  \node[box] (pairs) at (145,58) {同输入偏好对\\$(y_{\mathrm w},y_{\mathrm l})$};
  \node[box] (value) at (107,0) {价值模型\\前缀基线$V_{\bm\psi}(s_t)$};
  \node[box,fill=FigureLoss!10] (token) at (145,0) {逐token奖励\\$\widetilde r_t$};
  \node[box] (adv) at (145,-29) {回报与优势\\$\widehat A_t$};
  \node[box] (update) at (61,-29) {PPO式更新\\策略；另拟合价值};
  \node[box] (reference) at (107,-58) {固定参考策略概率\\$\pi_{\mathrm{ref}}(a_t\mid s_t)$};
  \draw[flow] (input) -- (policy);
  \draw[flow] (policy) -- node[lab,above] {生成} (response);
  \draw[flow] (response) -- (reward);
  \draw[flow] (pairs) -- node[lab,right] {偏好拟合} (reward);
  \draw[draw=BookMuted,line width=.8pt] (response.south) -- (107,15);
  \draw[flow] (107,15) -- node[lab,left] {前缀} (value.north);
  \draw[flow] (107,15) -- node[lab,above] {动作与$\mu$}
    (125,15) -- (125,0) -- (token.west);
  \draw[flow] (reward) -- node[lab,right] {序列分数} (token);
  \draw[flow] (token) -- (adv);
  \draw[flow] (value.south) -- (107,-14) -- (125,-14)
    -- (125,-29) -- (adv.west);
  \draw[flow] (adv.south) -- (145,-44)
    -- node[lab,above] {局部更新信号} (61,-44) -- (update.south);
  \draw[flow] (reference.east) -- node[lab,above] {参考概率}
    (167,-58) -- (167,0) -- (token.east);
  \draw[flow] (update) -- node[lab,left] {下一轮参数} (policy);
\end{tikzpicture}

  \caption{PPO式RLHF中的评价粒度。偏好标签训练完整回答评分器；旧策略逐token生成并记录采样概率，价值模型读取前缀提供基线。序列分数、旧策略概率与参考概率按式\eqref{eq:rl-reward-token-reward}构造逐token奖励，再经回报和优势估计进入策略更新；参考偏移不在更新器中重复计罚。图中连线表示数据和训练信号，不表示通过离散采样端到端反向传播。}
  \label{fig:rl-reward-rlhf}
\end{figure*}

逐token累加参考代价还会影响长短回答的权重。要把它与完整回答上的KL对应起来，可以使用序列概率分解所得的链式法则：
\[
  \operatorname{KL}(\pi(\cdot\mid q)\Vert\pi_{\mathrm{ref}}(\cdot\mid q))
  =\mathbb E_\pi\sum_{t<T}
    \log\frac{\pi(a_t\mid s_t)}{\pi_{\mathrm{ref}}(a_t\mid s_t)}.
\]
作为教学例，假设两个回答得分都为1，长度分别为2和4，各token的采样对数比恰好都是$0.1$。取$\beta=0.5$，扣除参考代价后的样本分数分别为$0.9$和$0.8$。差异来自累加项，并不说明奖励模型明确偏好短回答；真实对数比也未必同号或恒定。若将KL或策略损失除以长度，就会改变长短回答的相对权重，而不是只做无影响的数值缩放。选择哪种聚合应依据任务标准，不能将一种归一化当成通用答案。

偏好标签不必都由人提供。Constitutional AI按照明确原则生成批评和修订，结合监督修订与AI偏好，再由AI偏好学习奖励、优化策略\scite{bai2022constitutional}
给定评分函数后，更换评价者不改变优化式的形式，却改变评分依据和误差来源，这些条件将在第~\ref{sec:rl-reward-sources}节检查。现在先固定奖励$r$，考察完整回答的概率分布如何在奖励与参考偏移之间取舍。下面在有限输出集合上优化全部概率分布，以便精确求解；实际网络只允许选择其参数可表示的分布，范围更小。

\begin{theorem}[KL正则化奖励最大化的最优策略]
  \label{thm:rl-reward-kl}
  固定输入$q$，输出集合$\mathcal Y$有限，奖励$r(q,y)$有限，参考策略对每个$y\in\mathcal Y$均满足$\pi_{\mathrm{ref}}(y\mid q)>0$，且$\beta>0$。在所有概率分布上最大化
  \[
    F_q(\pi)=\mathbb E_{y\sim\pi(\cdot\mid q)}r(q,y)
     -\beta\operatorname{KL}(\pi(\cdot\mid q)\Vert
                                      \pi_{\mathrm{ref}}(\cdot\mid q))
  \]
  的唯一解为
  \begin{equation}
    \begin{aligned}
      \pi^*(y\mid q)&=
        \frac{\pi_{\mathrm{ref}}(y\mid q)\exp(r(q,y)/\beta)}{Z(q)},\\
      Z(q)&=\sum_{z\in\mathcal Y}\pi_{\mathrm{ref}}(z\mid q)
                                \exp(r(q,z)/\beta).
    \end{aligned}
    \label{eq:rl-reward-kl-optimum}
  \end{equation}
\end{theorem}

\begin{proof}
  有限输出、有限奖励和参考策略的正支持，使$0<Z(q)<\infty$，也使候选解$\pi^*$处处为正。为了比较任意$\pi$与这个候选解，由定义先将奖励写成
  $r(q,y)=\beta\log[\pi^*(y\mid q)/\pi_{\mathrm{ref}}(y\mid q)]
  +\beta\log Z(q)$。代入$F_q$后，奖励中的概率比与参考KL的概率比合并为
  \[
    F_q(\pi)=\beta\log Z(q)
       -\beta\sum_y\pi(y\mid q)\log\frac{\pi(y\mid q)}{\pi^*(y\mid q)}.
  \]
  第一项不依赖待优化的$\pi$，第二项是减去$\beta$倍的KL。由KL非负，只有$\pi=\pi^*$时第二项为零，故$\pi^*$唯一达到最大值。涉及零概率的项按$0\log0=0$连续延拓。
\end{proof}

最优解并不把全部概率放在最高分回答上，而是以参考概率为起点，按$\exp(r/\beta)$重新加权。假设两个输出的参考概率都是$1/2$，奖励为$(\log3,0)$，$\beta=1$，则$Z=2$，最优概率为$(3/4,1/4)$。若参考策略给某输出零概率，任何有限KL的策略也必须给它零概率，此时只能在参考支持内优化，不能沿用全输出空间的正支持结论。

这个精确解还提供了另一种偏好拟合方式。假定同一输入的偏好满足以$r$为效用的Bradley--Terry模型，便可把奖励差改写为最优策略相对参考策略的对数概率比差：
\begin{equation}
  \begin{aligned}
    r(q,y_{\mathrm w})-r(q,y_{\mathrm l})
    =\beta\left[
      \log\frac{\pi^*(y_{\mathrm w}\mid q)}
                    {\pi_{\mathrm{ref}}(y_{\mathrm w}\mid q)}
      -\log\frac{\pi^*(y_{\mathrm l}\mid q)}
                    {\pi_{\mathrm{ref}}(y_{\mathrm l}\mid q)}
    \right].
  \end{aligned}
  \label{eq:rl-reward-dpo-change}
\end{equation}
其中$y_{\mathrm w}$和$y_{\mathrm l}$分别表示胜出和落选回答。同一输入的两边都有$\beta\log Z(q)$，所以作差后消去；不同输入的归一化常数通常不能抵消。于是可以直接用策略概率比预测偏好，而不单独拟合奖励。以参数化策略替代$\pi^*$，将方括号内的量记为$z_{\bm{\theta}}(q,y_{\mathrm w},y_{\mathrm l})$，再最小化胜出标签的负对数似然，就得到\term{直接偏好优化}（direct preference optimization, DPO）：
\begin{equation}
  L_{\mathrm{DPO}}(\bm{\theta})
  =-\mathbb E_{(q,y_{\mathrm w},y_{\mathrm l})\sim\mathcal D_{\mathrm{pref}}}
         \log\sigma\!\left(\beta z_{\bm{\theta}}\right).
  \label{eq:rl-reward-dpo}
\end{equation}
期望在固定偏好记录$\mathcal D_{\mathrm{pref}}$上计算，每条记录含同一输入及其胜出、落选回答。一次更新计算两条完整回答在当前策略和固定参考下的对数概率，再按式~\eqref{eq:rl-reward-dpo}下降；这个步骤既不采集新轨迹，也不训练显式价值模型。单条记录的梯度为
$-\beta\sigma(-\beta z_{\bm{\theta}})
[\nabla\log\pi_{\bm{\theta}}(y_{\mathrm w}\mid q)
-\nabla\log\pi_{\bm{\theta}}(y_{\mathrm l}\mid q)]$。
沿负梯度更新，局部方向提高胜出回答相对落选回答的对数概率。沿用均匀参考的二输出例，策略logit差就是$z$；当$z=0$且$\beta=1$时，梯度为$-1/2$，以$0.2$为步长下降后$z=0.1$，胜出回答的概率约为$0.525$\scite{rafailov2023dpo}

显式奖励模型加PPO，要给新采样回答评分并更新策略；标准DPO只在固定偏好对上优化分类式损失，两者使用的数据和学习对象不同。上述换元解释了DPO与KL正则化最优策略的联系，却没有把DPO变成通常意义上的在线强化学习。实际模型类未必能表示$\pi^*$，有限记录未必覆盖所有回答，偏好也未必满足Bradley--Terry假设，所以不能由理想关系推出精确恢复最优策略。若在外层反复采样、标注再运行DPO，便形成迭代式偏好学习；在线访问和新目标信息来自外层协议，不来自分类损失本身。

% Outline: 5.2.5
\subsection{从过程与结果验证构造策略更新}
\label{sec:rl-reward-verification}

若任务有明确答案或可执行产物，我们还可以直接检查候选是否满足条件，而不逐对询问偏好。例如运行程序测试、按等价规则比较数学答案，或核对配送签收身份。这些检查可以重复采用同一标准，但它们只能评价实际检查到的条件，不能据此断言整个任务都已正确完成。

\begin{definition}[过程反馈与结果反馈]
  \label{def:rl-reward-verification}
  \term{结果反馈}（outcome feedback）评价行为结束后的产物或状态，例如答案是否等价、包裹是否正确签收。\term{过程反馈}（process feedback）评价中间状态、动作或步骤间关系，例如装载后身份是否一致、推导是否由前提合法推出。\term{可验证奖励}（verifiable reward）由明示验证规则产生；其可验证性依赖这条规则、输入前提和检查实现，不表示完整行为目标已经得到验证。
\end{definition}

过程与结果提供的证据也不能相互替代。推导可能使用非法除法却偶然得到正确答案，所以结果通过不足以证明过程有效；反过来，每步合法的变换也可能只是在反复改写题目，没有回答问题。表~\ref{tab:rl-reward-verifiers}据检查对象和前提比较验证器，自动检查是否可信仍要逐项判断。

\begin{table*}[htbp]
  \centering\small
  \begin{tabularx}{\linewidth}{@{}>{\raggedright\arraybackslash}p{21mm}*{4}{>{\raggedright\arraybackslash}X}@{}}
    \toprule
    检查器 & 可检查对象 & 检查前提 & 反馈粒度 & 主要漏检 \\
    \midrule
    规则验证器 & 格式、身份、答案等价 & 解析与等价规则符合标准 & 字段、步骤、结果 & 格式正确但内容错误；解析漏洞 \\
    执行器 & 代码、操作的实际结果 & 沙箱、测试输入、执行环境有效 & 工具调用或整次运行 & 未测试输入、超时与资源副作用 \\
    形式检查器 & 形式命题及证明 & 形式化忠实，公理和检查内核可信 & 推理步骤、完整证明 & 证明了错误形式化的命题 \\
    模型评判器 & 开放式内容与解释 & 训练覆盖、提示及标准适用 & 片段、回答 & 风格偏好、位置偏差、共同模型盲点 \\
    \bottomrule
  \end{tabularx}
  \caption{四种验证器实际检查的对象及漏检范围。执行测试只支持被测试输入，形式检查只支持所形式化的命题；两者都不能自行证明完整需求已被正确检查。}
  \label{tab:rl-reward-verifiers}
\end{table*}

拿到候选的验证得分后，可以比较同一输入下多个回答，为策略提供相对信号。\term{组相对策略优化}（group relative policy optimization, GRPO）以这组回答的得分为基准，给高于组均值的回答正信号，给低于均值的回答负信号，从而不必另拟合价值基线。沿用有限序列的未折扣评价，令旧策略$\mu$对$q$独立采样$m$个回答$y_i$，其验证或奖励模型得分记为$v_i$；这里$i$编号候选，不是$r_t$中的动作时刻。为中心化得分并统一组内尺度，定义
\begin{equation}
  \bar v=\frac1m\sum_{i=1}^m v_i,\qquad
  s_v=\sqrt{\frac1m\sum_{i=1}^m(v_i-\bar v)^2},\qquad
  \widehat A_i=\frac{v_i-\bar v}{s_v+\epsilon_{\mathrm{num}}},
  \label{eq:rl-reward-grpo-advantage}
\end{equation}
标准差使用分母$m$，$\epsilon_{\mathrm{num}}>0$防止分母为零。结果监督时，同一个$\widehat A_i$可以用于回答$i$的全部token，但它只表示该回答相对本组得分的高低，既不是逐token标签，也不是对状态价值函数的估计。

为了将这个组信号接入策略更新，令$T_i$为回答长度，$\varrho_{i,t}=\pi_{\bm{\theta}}(y_{i,t}\mid q,y_{i,<t})/
\mu(y_{i,t}\mid q,y_{i,<t})$为当前与旧策略的token概率比，$\ell_{i,t}^{\mathrm{clip}}$为式~\eqref{eq:rl-reward-ppo}中以$\widehat A_i$替代$\widehat A_t$的剪裁项。一个常见实现先在每个回答内平均token项，再在组内平均，并扣除参考偏移代价，最大化代理目标
\begin{equation}
  \widehat{\mathcal L}_{\mathrm{GRPO}}
  =\frac1m\sum_{i=1}^m\frac1{T_i}
      \sum_{t=0}^{T_i-1}
       \left(\ell_{i,t}^{\mathrm{clip}}-\beta k_{i,t}\right),
  \label{eq:rl-reward-grpo}
\end{equation}
其中$k_{i,t}$是所选参考KL估计或局部近似，常用形式为
$k=\exp(-\ell)+\ell-1$，$\ell=\log(\pi_{\bm{\theta}}/\pi_{\mathrm{ref}})$；
若当前策略和参考策略在同一合法token集合上都为正，且动作由$\pi_{\bm{\theta}}$采样，那么给定前缀有$\mathbb E_{\pi_{\bm{\theta}}}\exp(-\ell)=1$，从而$k$的期望等于前向KL。旧策略采样不自动满足这个分布条件；复用旧样本时若没有分布校正，样本平均就不能当成精确KL，更不能证明硬约束成立。长度平均、剪裁和采样KL都属于代理目标的选择，与无近似的期望奖励梯度需要区分。

\begin{example}[一组验证得分的更新信号]
  \label{ex:rl-reward-grpo}
  假设$m=4$，四个得分为$(1,1,0,0)$，则$\bar v=0.5$、$s_v=0.5$。忽略极小稳定项，相对信号为$(1,1,-1,-1)$。为计算一次更新，再设任务只有一步和两个动作，成功动作概率为$p=\sigma(z)$，当前$z=0$，不加KL惩罚，概率比处在未剪裁区。这组样本对logit的梯度为
  \[
    \frac14\{2(1)(1-p)+2(-1)(-p)\}=\frac12.
  \]
  用$0.2$为步长上升，得到$z'=0.1$，成功动作概率约为$0.525$。若整组全为1或全为0，中心化后的分子都为0；稳定项只能使计算有限，不能产生奖励梯度。若实现保留KL项，正则化仍可能产生梯度。
\end{example}

上述算例给出的是组内相对信号，不能直接当作无偏真实优势。即使不除以标准差，组均值含样本自身也会改变梯度期望。为了看清这一点，固定$q$和得分函数$v$，从$\pi$独立同分布采样，记$\bm g_i=\nabla\log\pi(y_i\mid q)$。利用$\mathbb E\bm g_i=0$，可得
\[
  \mathbb E\left[\frac1m\sum_i(v_i-\bar v)\bm g_i\right]
  =\left(1-\frac1m\right)\nabla\mathbb E_\pi v.
\]
缩放因子来自$\bar v$中与$\bm g_i$相关的自身得分，其余独立样本的交叉期望为零。再除以随机$s_v$，还会引入与组内成功比例有关的权重；二值奖励混合组中，较少见结果的标准化绝对值更大。改用不含自身的组均值可以消除这个特定缩放，但不能同时消除标准化、剪裁或复用旧样本造成的偏差。

组内比较只要求给出$v_i$，并未规定它必须来自规则。DeepSeekMath提出GRPO，并讨论结果奖励模型、过程奖励模型和迭代刷新，反馈并不限于规则判分\scite{rl5-shao2024deepseekmath}
DeepSeek-R1原始技术报告中的R1-Zero及R1推理训练使用正确性、格式等规则奖励；但R1完整多阶段流程还包含监督微调及一般数据上的模型奖励，不能用推理阶段的反馈概括全部训练\scite{deepseek2025r1}

基于可验证奖励的强化学习（reinforcement learning with verifiable rewards, RLVR）说明奖励由明示规则产生，GRPO则说明拿到分数后如何更新策略，两者不是同一个分类依据。验证器的假阳性、解析偏差、测试泄漏或长度与表达偏好，都可能被策略用来增加得分。增加同一输入的采样数可提供更多组内比较，却不会补齐未检查要求。在多轮工具任务中，还要明确评分针对动作格式、工具执行还是最终完成情况；将终局成功平均分给工具调用，并没有取得独立的动作正确性标签。更长交互的前沿问题见第~\ref{chap:rl-overview}章，语言模型实现由应用卷承接。

% Outline: 5.2.6
\subsection{反馈来源与目标可辨识性}
\label{sec:rl-reward-sources}

回到配送任务，已有示范可能足以复现动作，偏好记录可能足以预测路线胜负，验证器可能足以判断签收身份。要判断目标是否可辨识，必须先指定想要哪种结论，再检查产生同样反馈的候选解释是否都支持它。奖励即使不唯一，候选奖励也可能一致支持当前动作；偏好即使预测准确，也不能说明未被评价的要求已经学会。这正是前章可识别性在目标反馈上的应用。

同一对行为的胜负标签可以由人判断，也可以由模型判断。比较这两种来源，主要要检查判断依据、成本和误差是否相同，而不是仅比较记录格式。形式相同不意味着评价质量相同。

\term{基于AI反馈的强化学习}（reinforcement learning from AI feedback, RLAIF）由AI承担部分评价，再用所得反馈学习策略。\term{原则指导反馈}（constitutional feedback）则以明确原则指导评价或修订，用这些原则产生批评、修订示范或偏好。前者说明评价来源，后者说明判断依据；原则不是与示范、偏好互斥的记录类型，其选择仍须由外部目标约定。

\begin{table*}[t]
  \centering\small
  \begin{tabularx}{\linewidth}{@{}>{\raggedright\arraybackslash}p{22mm}*{5}{>{\raggedright\arraybackslash}X}@{}}
    \toprule
    来源或协议 & 判断依据 & 相关偏差 & 可追溯性 & 查询成本 & 外部校验 \\
    \midrule
    人类反馈 & 任务知识、个人偏好、指南 & 疲劳、群体差异、顺序效应 & 需保留指南与标注群体 & 人工时间 & 交叉复核、真实结果 \\
    RLAIF & AI知识与提示中的标准 & 共享训练偏差、风格偏好 & 需记录模型和提示版本 & 推理及抽查成本 & 独立人评、规则或执行 \\
    原则指导反馈 & 显式原则及其解释 & 原则遗漏、冲突解释 & 原则可见不等于执行透明 & 原则维护与评价成本 & 逐原则反例测试 \\
    模型评判器 & 学得评分或比较函数 & 分布外失准、位置与长度偏差 & 需记录训练数据和模型版本 & 批量打分与刷新 & 留出来源、实物或程序检查 \\
    \bottomrule
  \end{tabularx}
  \caption{反馈来源及协议的判断依据、成本与校验方式。RLAIF可使用原则指导的模型评判器，因此同一系统可以同时属于多行，不能按互斥类别解读。}
  \label{tab:rl-reward-sources}
\end{table*}

假设人工标注员和模型都只看终点照片，两者都判断“已送达”，却都看不到包裹内部破损。增加来源一致性没有增加内部状态的信息，因此不能推出包裹完好。来源分歧也未必是噪声：一位收件人看重速度，另一位看重省电，可能确实采用不同取舍，不应直接删除其中一方的反馈。

表~\ref{tab:rl-reward-sources}列出的外部校验，也应带上输入和任务范围。评价器只在信息完整的短距离记录上检查过，就不能直接用于长路线、身份缺失或新包裹，并声称同样可靠。现有证据只支持已检验策略在这些输入和标准下的结论；要判断多个奖励解释是否会造成实际失配，还要比较未见行为上的排序。

% Outline: 5.3
\section{奖励归因受限下的强化学习}
\label{sec:rl-reward-credit}

假设我们已约定配送成功的含义，也能取得相应终局反馈，仍要决定怎样用这个结果更新核对、装载等早期动作。本节暂时固定目标和反馈语义，比较结果回溯、过程判断、奖励变换与辅助目标。这些办法都能为局部决策提供信号，但有的增加信息，有的只改变记账或训练方式，保持原目标所需的条件也不同。

% Outline: 5.3.1
\subsection{基于长期结果回溯构造局部学习信号}
\label{sec:rl-reward-redistribution}

在四步配送中，早期核对决定是否拿对包裹，中间搬运可能只是执行既定方案。REINFORCE或多步回报能将终局成功用于核对动作的更新，但搬运和提示播放也可能收到同一个成功回报。前章的基线和自举分别减少共同波动、缩短传播路径，并没有直接测量“改变这个动作会使结果变化多少”。这里回用第~\ref{chap:rl-learning}章的估计方法，进一步检查局部信号所表达的贡献。

\term{奖励重分配}（reward redistribution）将同一结果重新记到轨迹上的不同位置，使结果与关键事件的时间联系更短。要判断这种记账是否仍优化原目标，必须按原来的折扣权重比较整条轨迹的新旧奖励，而不能只比较未折扣总和。若要求逐轨迹保持回报，应满足
\begin{equation}
  \sum_{t=0}^{H-1}\gamma^t\widetilde r_t
    =\sum_{t=0}^{H-1}\gamma^t r_t.
  \label{eq:rl-reward-redistribution}
\end{equation}
逐轨迹相等足以保持每个策略的期望回报；也可以使用较弱条件，但需要期望相等对全部待比较策略成立。另要检查信号何时可得。依赖整段历史的新奖励通常不是原状态上的马尔可夫奖励，需要相应表示；若它依赖记账时刻之后才发生的结果，只加入过去历史仍不能使反馈当时可得。

例如$H=4$、$\gamma=0.9$，原奖励$(0,0,0,1)$对应初始回报$0.9^3=0.729$。若全部移至$t=0$，新奖励应为$(0.729,0,0,0)$；若移至$t=1$，则应为$(0,0.81,0,0)$，两者按折扣计算仍为$0.729$。直接移成$(1,0,0,0)$只保留未折扣总和，没有保留原目标。若轨迹长度可变，同样忽略折扣还会改变快慢路径的排序。

保持总回报也不意味着可以直接用新奖励计算reward-to-go，即从当前动作起的后续奖励回报，包含该动作的奖励。假设终局成败取决于最后一个动作，却将全部结果事后记到$t=0$，末步的局部回报就变为0，用它更新末步策略会丢掉该动作的影响。第~\ref{chap:rl-learning}章之所以能删去过去奖励，是因为过去结果不受当前动作影响；事后前移的信号不满足这个因果条件。将不变的完整轨迹回报与整条轨迹的对数概率梯度相乘，仍保留原目标梯度。要改用局部更新，还须另外证明重分配与估计器的配套条件。

一种重分配依据是前缀所提供的结果预测信息。令$m_t$预测看到前缀后的最终回报，将$m_t-m_{t-1}$视为新增预测信息；取$m_{-1}=0$且$m_{H-1}=G_0$，相邻差求和恰为$G_0$。若将其记为折扣局部奖励，且$0<\gamma\leq1$，可取$\widetilde r_t=\gamma^{-t}(m_t-m_{t-1})$以抵消记账位置的折扣。末端预测不等于真实回报时必须补残差，否则总回报也改变了。实际学习还需控制预测误差和$\gamma^{-t}$的放大；RUDDER用回报预测与贡献分析构造信号，其机制入口已在第~\ref{chap:rl-learning}章介绍。

预测信息仍不等于干预贡献。假设正确核对后屏幕总会闪光，日志中闪光与成功高度相关；但关闭闪光不会改变交付，拿错包裹才会导致失败。归因模型只依据相关性时，可能把奖励分给闪光。检验应改变无关标记、受控比较关键动作，并观察实际交付结果；仅有回报预测误差下降不足以证明贡献判断正确。

% Outline: 5.3.2
\subsection{基于过程反馈构造局部学习信号}
\label{sec:rl-reward-process}

若终局日志无法区分核对错误与中途换包，可以新增装载前的身份检查。这个检查提供新的中间证据，而不是将原总分分摊到各步。需要先明确它评价什么：读码动作、装载后的身份关系和到达时的包裹状态对应不同决策，只有步骤边界和被检关系清楚，才能确定反馈监督的对象。

图~\ref{fig:rl-reward-process}把人工标注、规则检查、执行记录和模型步骤评价放到同一轨迹中，标明各自覆盖的步骤。例如执行记录只证明包裹放入箱内，不证明身份正确；模型觉得搬运平稳，也不能替代破损检查。推理任务中的逐步骤验证是过程监督的一个代表实例，同样需要明确标签标准及标注范围\scite{lightman2023verify}

即使各处检查都通过，局部正确也可能没有组成全局成功。机器人每次搬运都遵守速度限制、每次放置都平稳，却只在两个中间站往返，没有交给收件人。过程条件没有覆盖最终交付；若各次身份检查只看单个字段，还可能遗漏跨步骤身份一致性。因此，将局部分数求和之前要解释它与整体要求的关系，终局检查不能仅因步骤有标准就取消。

过程反馈也增加获取和维护成本：标注量通常随步骤数增长，不同标注者还可能采用不同的步骤切分。模型反馈若有系统偏差，会在多个位置重复传入更新。可以比较只补关键步骤与标注全部步骤所需成本，统一标准版本，再检查独立交付结果是否改善。若发现检查标准遗漏了要求，就需要修订目标标准，不能只调整局部信号权重。

\begin{figure*}[htbp]
  \centering
  % Each feedback evaluates a stated step, not every preceding action.
\begin{tikzpicture}[
  x=1mm,y=1mm,font=\figurefont\fontsize{8}{10}\selectfont,
  box/.style={draw=BookRule,fill=BookPaper,text width=23mm,
    align=center,minimum height=10mm,inner sep=2mm,line width=.7pt},
  flow/.style={-{Stealth[length=2mm]},draw=BookMuted,line width=.8pt},
  signal/.style={flow,draw=BookAmber,dashed},
  lab/.style={fill=white,inner sep=.5mm,font=\figurefont\fontsize{7}{9}\selectfont}
]
  \node[box] (a0) at (15,23) {核对身份\\$a_0$};
  \node[box] (a1) at (58,23) {装载\\$a_1$};
  \node[box] (a2) at (101,23) {搬运（含选路）\\$a_2$};
  \node[box] (a3) at (144,23) {交付\\$a_3$};
  \draw[flow] (a0) -- (a1);
  \draw[flow] (a1) -- (a2);
  \draw[flow] (a2) -- (a3);
  \node[box,fill=FigureInput!10] (f0) at (15,-1) {人工标注\\读码动作};
  \node[box,fill=FigureInput!10] (f1) at (58,-1) {规则检查\\装载身份一致};
  \node[box,fill=FigureInput!10] (f2) at (101,-1) {模型步骤评价\\只评平稳程度};
  \node[box,fill=FigureInput!10] (f3) at (144,-1) {执行记录\\交付结果};
  \foreach \i in {0,1,2,3}{
    \draw[flow] (a\i) -- node[lab,right] {被评对象} (f\i);
    \node[box,fill=FigureModel!10] (u\i) at (\i*43+15,-26) {相应步骤的\\策略更新信号};
    \draw[signal] (f\i) -- node[lab,right] {反馈} (u\i);
  }
  \node[anchor=north,align=center] at (80,-38)
    {早期选路没有直接标签，仍靠结果回溯；终局身份与完好要求另行检验};
\end{tikzpicture}

  \caption{过程反馈缩短部分决策到评价的路径。上方为按时间执行的配送步骤，下方为评价来源；由反馈指向更新对象的虚线表示学习信号。未标注步骤仍依赖终局回报传播，局部通过不取代最终身份和完好检查。}
  \label{fig:rl-reward-process}
\end{figure*}

% Outline: 5.3.3
\subsection{基于奖励变换构造局部学习信号}
\label{sec:rl-reward-shaping}

即使不新增检查，也可以依据已有目标构造即时信号。例如希望鼓励机器人接近终点，但直接给接近动作加分可能使它反复离开再靠近。要避免附加分数任意改变累计目标，可以采用\term{潜势型奖励塑形}（potential-based reward shaping）：为每个状态设一个评分，按与任务相同的折扣计算相邻评分差，使沿途附加项能够相消。

\begin{definition}[潜势型塑形]
  \label{def:rl-reward-shaping}
  给定与任务相同的折扣$\gamma$及状态势函数$\Phi:\mathcal S\to\mathbb R$，定义
  \[
    F(s,a,s')=\gamma\Phi(s')-\Phi(s),\qquad
    r'_t=r_t+\gamma\Phi(s_{t+1})-\Phi(s_t).
  \]
  $\Phi$是人为指定或从数据学习的状态评分，称为势函数；附加奖励必须采用这个折扣差形式，不能用任意启发式分数替代。
\end{definition}

为判断附加信号是否仍评价原任务，我们比较塑形前后的完整折扣回报。下面利用相邻势项相消，证明在给定条件下两者只差一个不依赖策略的期望常数；这是策略不变奖励变换的经典结论\scite{rl5-ng1999shaping}

\begin{theorem}[无限折扣时域的塑形不变性]
  \label{thm:rl-reward-shaping}
  设$0\leq\gamma<1$，$\Phi$有界，原折扣回报存在且可积，各待比较策略使用相同初始分布$\rho_0$。则塑形后的回报满足
  \[
    G'_0=G_0-\Phi(s_0),\qquad
    J'(\pi)=J(\pi)-\mathbb E_{\rho_0}\Phi(s_0).
  \]
  因而所有策略在标准期望折扣回报下的排序不变。
\end{theorem}

\begin{proof}
  先取长度为$H$的有限前缀，以便保留末端边界项。将塑形奖励代入折扣和，相邻势项逐项相消：
  \begin{align*}
    \sum_{t=0}^{H-1}\gamma^t r'_t
    &=\sum_{t=0}^{H-1}\gamma^t r_t
       +\sum_{t=0}^{H-1}
          \bigl[\gamma^{t+1}\Phi(s_{t+1})
                          -\gamma^t\Phi(s_t)\bigr]\\
    &=\sum_{t=0}^{H-1}\gamma^t r_t
       -\Phi(s_0)+\gamma^H\Phi(s_H).
  \end{align*}
  由$|\Phi(s)|\leq M$，末端项绝对值至多为$\gamma^H M$，当$H$趋于无穷时消失；若$\gamma=0$，从$H\geq1$起该项就为零。因此每条轨迹满足$G'_0=G_0-\Phi(s_0)$。势的有界性保证差值可积，取期望得到只依赖共同初始分布的偏移。任意两策略减去同一个期望数值，其回报差和排序均不变。
\end{proof}

保持期望排序不要求$\Phi$等于精确价值。即使势函数误将慢路线评为接近终点，只要有界且按同一折扣差加入，定理仍成立；但错误势值可能妨碍学习，使速度或有限样本回报变差。若初始状态随机，$-\Phi(s_0)$也是随机偏移，可能改变方差及条件风险价值（conditional value-at-risk, CVaR）这类尾部评价。以损失$-G_0$定义CVaR时，它评价上尾的大损失，对应回报较差的结果。第~\ref{chap:rl-trustworthy}章的风险目标需要另查不变性，不能继承这里的期望结论。

有限任务没有无穷尾项可供消去，必须保留终端势。允许势随时刻变化，采用
$r'_t=r_t+\gamma\Phi_{t+1}(s_{t+1})-\Phi_t(s_t)$。
沿相同方式展开，可得
\begin{equation}
  \sum_{t=0}^{H-1}\gamma^t r'_t
  =\sum_{t=0}^{H-1}\gamma^t r_t
      -\Phi_0(s_0)+\gamma^H\Phi_H(s_H).
  \label{eq:rl-reward-shaping-terminal}
\end{equation}
这个恒等式也适用于有限时域的$\gamma=1$。令$\Phi_H=0$足以消去末端项；若$H$固定，终端势为所有策略共享的常数也能保持排序。$\gamma=0$时末端项不影响目标，但其他折扣下必须显式处理边界。

自然终止和采样时间截断对应不同边界。自然终止后不再记账时，应将终止势设为零，或显式抵消$\gamma^T\Phi_T(s_T)$；若$T$可变且$\gamma<1$，即便终端势都是同一非零常数，折扣后的末端项也会随长度变化。若采用吸收状态并无限继续记塑形奖励，尾项可以继续相消，但不能只记进入吸收状态的一步，省略随后塑形项。训练批次截断则可能尚未结束任务，需要保留原续接价值和对应塑形价值$V'(s)=V(s)-\Phi(s)$；若不保留，就要明示改求有限时域目标。未计算续接价值不意味着它为零。

\begin{example}[遗漏终端势会改变排序]
  \label{ex:rl-reward-terminal}
  假设$H=1$、$\gamma=0.9$，从同一起点选择到终点甲或乙，原奖励都是0，且$\Phi_0(s_0)=0$。若$\Phi_1(\text{甲})=1$、$\Phi_1(\text{乙})=0$，到达后立即停止记账，新回报分别为$0.9$和0，原本并列的路径变成甲更优。将两个终端势都设为0则消除了这一差异；在含中间步骤的任务中，即使中间势不准确，只要满足相应定理条件，仍可保持期望排序。
\end{example}

现在可以具体区分这些保持条件。潜势塑形使共同起点的回报只差固定势值，重分配可以要求逐轨迹折扣回报完全相同，一般辅助奖励则不自动具有这两种性质。即使附加信号更密集、代数上保持了目标，也没有证明它等于动作的因果贡献。

% Outline: 5.3.4
\subsection{基于辅助目标构造局部学习信号}
\label{sec:rl-reward-auxiliary}

让机器人预测下一时刻位置，可以用预测误差训练表示，使其保留动力学信息；为每次到达装载区加分，则直接改变策略优化的回报。两者都可提供辅助信号，但前者作用于表示训练，后者作用于行为计分，既不都是新评分，也不都是原回报的等价变换。

为说明更新位置，设共享表示参数为$\bm{\eta}$，任务的策略或价值损失为$L_{\mathrm{task}}$，下一状态预测损失为$L_{\mathrm{aux}}$。以$\lambda$设置辅助损失权重、$\alpha$为步长，联合更新为
$\bm{\eta}^{(k+1)}=\bm{\eta}^{(k)}
-\alpha\nabla_{\bm{\eta}}(L_{\mathrm{task}}+\lambda L_{\mathrm{aux}})$。
辅助预测梯度改变共享表示，最终行为仍按原任务回报评价。若改为给动作增加辅助奖励$b_t$，采用$r'_t=r_t+\lambda b_t$，则优化目标变为
$J'(\pi)=J(\pi)+\lambda\mathbb E_\pi\sum_t\gamma^t b_t$，
除非$b_t$满足额外不变条件，这个新目标不再等于原任务。图~\ref{fig:rl-reward-auxiliary}将两条路径分别接到共享表示和回报，显示它们改变的对象。

\begin{figure*}[htbp]
  \centering
  % Two interventions: representation gradients versus reward aggregation.
\begin{tikzpicture}[
  x=1mm,y=1mm,font=\figurefont\fontsize{8}{10}\selectfont,
  box/.style={draw=BookRule,fill=BookPaper,text width=25mm,
    align=center,minimum height=10mm,inner sep=2mm,line width=.7pt},
  flow/.style={-{Stealth[length=2mm]},draw=BookMuted,line width=.8pt},
  grad/.style={flow,draw=BookAmber,dashed},
  lab/.style={fill=white,inner sep=.7mm,font=\figurefont\fontsize{7}{9}\selectfont}
]
  \node[anchor=west] at (0,45) {A\quad 辅助预测损失};
  \node[box,fill=FigureInput!10] (obs) at (35,29) {轨迹观测};
  \node[box,fill=FigureModel!10] (rep) at (35,7) {共享表示\\参数$\bm\eta$};
  \node[box,fill=FigureLoss!10] (task) at (15,-20) {任务损失\\$L_{\mathrm{task}}$};
  \node[box,fill=FigureLoss!10] (aux) at (57,-20) {辅助预测损失\\$L_{\mathrm{aux}}$};
  \draw[flow] (obs) -- (rep);
  \draw[flow] (rep.south) -- ++(0,-4) -| (task.north);
  \draw[flow] (rep.south) -- ++(0,-4) -| (aux.north);
  \draw[grad] (task.west) -- ++(-6,0) |- node[lab,near end,above] {梯度} (rep.west);
  \draw[grad] (aux.east) -- ++(6,0) |- node[lab,near end,above] {梯度} (rep.east);
  \node[anchor=west] at (88,45) {B\quad 加入辅助奖励};
  \node[box,fill=FigureInput!10] (r) at (101,29) {原奖励$r_t$};
  \node[box,fill=FigureInput!10] (b) at (145,29) {辅助奖励$b_t$};
  \node[box] (sum) at (123,7) {新回报\\$\sum_t\gamma^t(r_t+\lambda b_t)$};
  \node[box,fill=FigureModel!10] (pi) at (123,-20) {策略更新};
  \draw[flow] (r.south) -- ++(0,-5) -| (sum.north);
  \draw[flow] (b.south) -- ++(0,-5) -| (sum.north);
  \draw[flow] (sum) -- node[lab,right] {回报／优势} (pi);
  \node[anchor=north,align=center] at (81,-36)
    {共同评价：原任务交付结果与成本；辅助表现另报};
\end{tikzpicture}

  \caption{辅助预测损失与辅助奖励的干预位置。左侧虚线为损失对共享表示的梯度；右侧辅助奖励先进入回报，再产生策略更新信号。两种训练的最终行为都应按未掺入辅助分数的原任务标准评价。}
  \label{fig:rl-reward-auxiliary}
\end{figure*}

另外两种训练安排也需区分。训练课程改变任务分布，例如先练短路线，再增加距离；人为子目标改变决策组织，例如先到装载区，再去收件处。二者都不必增加奖励，也不等于潜势塑形。时间抽象回用第~\ref{chap:rl-learning}章；目标重标记及事后经验回放（hindsight experience replay, HER），即将轨迹按实际达到的目标重新使用，由第~\ref{chap:rl-transfer}章展开。

若进入装载区就给辅助分，机器人可能反复进出而不交付，因为计分没有限制重复次数，也未要求主任务进展。调整$\lambda$、切换课程或移除辅助信号后，应分别检查原始交付成功、辅助预测误差和子目标完成情况，不能用辅助分的上涨替代主任务收益。表示预测更准确也不证明动作贡献判断正确；共享参数上的辅助梯度还可能与任务梯度冲突。

% Outline: 5.3.5
\subsection{奖励归因方法的误差与适用边界}
\label{sec:rl-reward-credit-boundaries}

同一终局结果可以回溯记账，可以补充过程检查，也可以与塑形或辅助目标结合。要在这些办法间选择，我们需要比较信息获取成本、局部信号的含义和原目标的保持条件。表~\ref{tab:rl-reward-credit-methods}按这几个共同对象对照四条路线。

\begin{table*}[htbp]
  \centering\small
  \begin{tabularx}{\linewidth}{@{}>{\raggedright\arraybackslash}p{23mm}*{5}{>{\raggedright\arraybackslash}X}@{}}
    \toprule
    路线 & 额外反馈成本 & 所需轨迹信息 & 局部信号 & 原目标保持条件 & 主要失效 \\
    \midrule
    长期结果回溯 & 可无新标签；需拟合 & 前缀、终局结果 & 回报预测变化或重分配 & 同一折扣下的回报相等，或全策略期望相等 & 相关标记代替贡献、残差遗漏 \\
    过程反馈 & 关键步骤标注或检查 & 被评步骤及上下文 & 局部评分或通过标记 & 若加入回报，须另证聚合后仍保持原策略排序 & 局部正确不能组合为成功 \\
    奖励变换 & 势函数构造成本 & 相邻状态、终止类型 & 折扣势差 & 有界势、共同初始分布、正确终端项 & 截断当终止、折扣不一致 \\
    辅助目标 & 预测标签、额外任务 & 表示或子任务信息 & 辅助梯度、额外奖励 & 预测损失不改评价；额外奖励另证等价 & 梯度冲突、子目标循环 \\
    \bottomrule
  \end{tabularx}
  \caption{奖励归因路线的适用条件。过程反馈和一般辅助奖励没有自动保持原目标的代数保证。}
  \label{tab:rl-reward-credit-methods}
\end{table*}

即使采用同一条路线，误差也可能来自不同环节。归因模型错误把已经明确的成功结果分给无关动作；有限样本误差使现有证据不足以区分贡献；反馈标准错误则使“成功”本身遗漏破损要求。增加样本可能减轻第二种误差，但不自动修正第一种模型假设，也不会让错误标准变完整。

可以在同一配送任务中固定策略架构与预算，分别检查这些缺口。移除身份过程标签后看身份错误是否重新出现，用来判断中间信息的作用；打乱不影响环境的屏幕闪光，再看局部信号和策略是否改变，用来检测相关性捷径。保持原奖励不变，核查自然终止和时间截断的势值处理，则可排除记账意外改变目标。各项检查都要记录原始交付结果及错误类型，不能只比较掺入辅助分的奖励曲线。通用信用传播仍由前章工具承担，本节补充其信号来源和语义；标准本身的缺口需要转向目标诊断。

% Outline: 5.4
\section{奖励与行为目标不匹配}
\label{sec:rl-reward-mismatch}

配送策略可能得到许多反馈，也能将反馈传给早期动作，却仍以高分完成错误交付。此时需要检查计分标准与行为要求的关系。本节按诊断的实际依赖展开：先写清要求和证据，再检测偏差、判断根因，据此修订目标或估计方法。修订后还要用新行为重新检测，更换奖励名称并没有解决偏差。

% Outline: 5.4.1
\subsection{明确行为目标和代理奖励}
\label{sec:rl-reward-proxy}

配送要求包括交给正确收件人、保持包裹完好，以及满足约定的代价条件，训练代理却可能只给终点图像打分。继续用同一个评分器评价策略，就仍看不到内部破损。表~\ref{tab:rl-reward-requirements}将每项要求与训练编码、独立证据和未覆盖条件对应；这种逐项检查能显示平均总分可能隐藏的遗漏。

\begin{table*}[htbp]
  \centering\small
  \begin{tabularx}{\linewidth}{@{}*{4}{>{\raggedright\arraybackslash}X}@{}}
    \toprule
    行为要求 & 训练编码 & 独立评价 & 未覆盖条件 \\
    \midrule
    交给正确收件人 & 终点照片评分 & 签收身份与包裹编号双重核对 & 照片不含身份时仍缺证据 \\
    包裹完好 & 平稳搬运过程分 & 独立开箱检查或有效传感记录 & 隐性破损与传感盲区 \\
    能耗可接受 & 电量成本或偏好权重 & 校准电表、同负载路线比较 & 新载荷、不同电池状态 \\
    按时交付 & 终局成功或到达时间分 & 固定截止时间下的完成记录 & 长尾延迟与未完成轨迹 \\
    \bottomrule
  \end{tabularx}
  \caption{先建立行为要求到证据的对应，再选择数值聚合。某项无法观测意味着结论受限，不意味着它可以从目标中消失。}
  \label{tab:rl-reward-requirements}
\end{table*}

完整记录要求也不意味着所有要求都能同时改善。例如高收益配送可能更耗电，单一总分已经替决策者选定了取舍。向量奖励保留各项目标的数值，再另行指定怎样比较它们，使权衡不被藏在总分中。

\begin{definition}[向量奖励、标量化与Pareto最优]
  \label{def:rl-reward-pareto}
  令$\bm r_t=(r_{t,1},\ldots,r_{t,d})^\top$记录$d$个目标，期望回报向量为
  $\bm J(\pi)=\mathbb E_\pi\sum_t\gamma^t\bm r_t$。
  \term{标量化}（scalarization）用给定函数把多个目标转成一个数值，例如以$\bm w^\top\bm J(\pi)$比较策略，其中$\bm w\geq0$规定各目标的权重。
  在指定策略类$\Pi$中，若不存在$\pi'\in\Pi$使每个目标都满足$J_j(\pi')\geq J_j(\pi)$且至少一个严格大于，则称$\pi$为\term{Pareto最优}（Pareto optimal）。也就是说，在这组候选中，改善某个目标必须牺牲至少一个其他目标。所有分量约定为越大越好，能耗可用负能耗表示。
\end{definition}

图~\ref{fig:rl-reward-pareto}采用四个教学候选：甲$(2,4)$、乙$(4,7)$、丙$(6,8)$、丁$(5,5)$，横、纵坐标分别为能耗和收益，所以向左上改善。乙比丁收益高且能耗低，因而支配丁；甲、乙、丙在这四个候选中没有被支配。若按“收益减能耗”比较，三者得分为2、3、2，就选择乙。改变策略类，例如允许随机混合，会改变可达集合及候选前沿。若标量化非线性，还需规定先取回报期望再变换，还是先变换随机回报再取期望，两种次序一般不等价。多目标策略集的求解可作延伸阅读\scite{rl5-vanmoffaert2014pareto}

\begin{figure}[htbp]
  \centering
  % Constructed candidate policies: (energy, benefit).
\begin{tikzpicture}
  \begin{axis}[
    width=.96\linewidth,height=62mm,
    xmin=0,xmax=7,ymin=0,ymax=10,
    xlabel={能耗（任意单位，越小越好）},
    ylabel={收益（任意单位）},
    xtick={0,2,4,6},ytick={0,2,4,6,8,10},
    axis lines=left,axis line style={BookMuted,line width=.7pt},
    tick style={BookMuted},grid=major,grid style={BookRule,line width=.5pt},
    font=\figurefont\fontsize{8}{10}\selectfont,
    label style={font=\figurefont\fontsize{8}{10}\selectfont},
    tick label style={font=\figurefont\fontsize{7}{9}\selectfont},
    clip=false
  ]
    \addplot[BookTeal,dashed,line width=.9pt] coordinates {(2,4) (4,7) (6,8)};
    \addplot[only marks,mark=*,mark size=2.5pt,BookTeal]
      coordinates {(2,4) (4,7) (6,8)};
    \addplot[only marks,mark=square,mark size=2.7pt,BookMuted,line width=.9pt]
      coordinates {(5,5)};
    \node[anchor=east] at (axis cs:1.9,4) {甲};
    \node[anchor=south east] at (axis cs:3.9,7) {乙};
    \node[anchor=south] at (axis cs:6,8.1) {丙};
    \node[anchor=west] at (axis cs:5.15,5) {丁：被支配};
    \draw[-{Stealth[length=1.8mm]},BookMuted,line width=.8pt]
      (axis cs:4.9,5.2) -- (axis cs:4.15,6.7);
  \end{axis}
\end{tikzpicture}

  \caption{收益越大、能耗越小越好。实心点是在四个教学候选策略中的非支配点，空心方块丁被乙支配；虚线只帮助识别候选前沿，不断言连线上的策略全部可实现。}
  \label{fig:rl-reward-pareto}
\end{figure}

据此可以区分几种看似相近的安排。预先给定偏好后求一个策略，与保留多个非支配策略供以后选择，求解的输出不同；多个来源都估计“完好交付”，则仍是同一目标的多份证据。若规定能耗必须低于上限，能耗成为不可任意用收益补偿的约束，而不只是加权分数。约束求解与执行保护由第~\ref{chap:rl-trustworthy}章展开。

% Outline: 5.4.2
\subsection{检测行为与代理奖励的偏差}
\label{sec:rl-reward-detection}

用配送代理给行为打分以前，要先检查代理是否遗漏了已知要求；策略更新以后，还要检查新行为是否超出原反馈覆盖。两种检查都需要在优化前固定独立指标、留出反馈来源和必要情形。例如按路线长度、身份信息完整性与包裹类型分层，可以主动检查训练成功轨迹中少见的情况，只从常见成功轨迹抽样则可能继续漏掉它们。

图~\ref{fig:rl-reward-overoptimization}用解析曲线演示一种待检测现象：代理分数一直上升，按独立标准衡量的行为可能先改善再恶化。两条曲线的差别提醒我们同时查看训练代理和任务证据，不说明所有训练都会遵循这个形状。真实比较还需报告指标口径、采样预算、误差范围与模型选择过程。Gao等研究了奖励模型过度优化的经验现象，本图仍是教学示意，不是该论文的测量结果\scite{gao2022overoptimization}

不同测试可以针对不同缺口。反例测试检查已有怀疑，例如照片看似送达但身份错误；压力测试增加路线长度、载荷或重复次数，考察长时交互及极端输入。对抗测试主动寻找骗过代理的行为，例如展示别人的签收码；人工审查则分析难以形式化的失败。每次都应记录具体输入、行为、分数和验证器版本，使失败可以复查，不能只留下汇总总分。

如果依据留出结果调整奖励或挑选策略，这批证据就参与了开发。下一轮要补充新的独立检查，或明确报告测试复用范围，不能继续把旧结果当作完全未见数据上的验证。即使没有发现偏差，结论也只适用于已测情形，不能证明所有未见行为均满足要求。

\begin{figure*}[htbp]
  \centering
  % Analytic illustration only: p(x)=.2+.7*x, b(x)=.2+1.2*x-x^2.
\begin{tikzpicture}
  \begin{axis}[
    name=proxy,width=76mm,height=61mm,
    xmin=0,xmax=1,ymin=0,ymax=1,
    xlabel={优化强度$x$},ylabel={分数（任意单位）},
    title={A\quad 代理奖励},
    xtick={0,.5,1},ytick={0,.5,1},
    axis lines=left,axis line style={BookMuted,line width=.7pt},
    tick style={BookMuted},grid=major,grid style={BookRule,line width=.5pt},
    font=\figurefont\fontsize{8}{10}\selectfont,
    tick label style={font=\figurefont\fontsize{7}{9}\selectfont},
    domain=0:1,samples=81
  ]
    \addplot[BookBlue,line width=1.1pt] {.2+.7*x};
  \end{axis}
  \begin{axis}[
    at={(proxy.outer east)},anchor=outer west,xshift=8mm,
    width=76mm,height=61mm,xmin=0,xmax=1,ymin=0,ymax=1,
    xlabel={优化强度$x$},ylabel={分数（任意单位）},
    title={B\quad 独立任务表现},
    xtick={0,.5,1},ytick={0,.5,1},
    axis lines=left,axis line style={BookMuted,line width=.7pt},
    tick style={BookMuted},grid=major,grid style={BookRule,line width=.5pt},
    font=\figurefont\fontsize{8}{10}\selectfont,
    tick label style={font=\figurefont\fontsize{7}{9}\selectfont},
    domain=0:1,samples=81
  ]
    \addplot[BookAmber,dashed,line width=1.1pt] {.2+1.2*x-x^2};
    \node[anchor=south] at (axis cs:.6,.63) {先升后降};
  \end{axis}
\end{tikzpicture}

  \caption{代理与独立表现可能分离的双面板示意。横轴$x\in[0,1]$是归一化优化强度，曲线分别由$p(x)=0.2+0.7x$与$b(x)=0.2+1.2x-x^2$构造，纵轴为任意单位；不是实验结果，也不表示固定的普遍动力学。}
  \label{fig:rl-reward-overoptimization}
\end{figure*}

% Outline: 5.4.3
\subsection{区分奖励歧义与行为目标偏差}
\label{sec:rl-reward-ambiguity}

两个奖励都解释配送示范时，我们要继续问它们在哪些行为上等价。第~\ref{sec:rl-reward-shaping}节的势差在边界条件满足时保持所有策略的标准期望回报排序，只改变表示；另一些奖励却只在已见示范上相容，在未见路径上作出不同选择。后者限制目标外推，但要认定当前行为失配，仍需比较该行为与独立要求。

\begin{example}[示范相容不等于新路径等价]
  \label{ex:rl-reward-ambiguity}
  假设训练环境只允许专家路径$E$和路径$B$，两个终局奖励都给它们$(2,0)$，于是都支持专家选择$E$。测试环境关闭$E$，另开放此前不可行且未经专家比较的$C,D$。$R_1$给$(C,D)$的分数为$(1,-1)$，$R_2$给$(-1,1)$，而两者仍给$B$记0分。因此，对训练环境中每一条可行路径，两种奖励给出的分数相同；不同路径之间的分数并不相同。测试时两种奖励却分别支持选择$C$和$D$。即使再重复收集一万次训练选择，新增记录仍不涉及$C,D$，不能区分两种目标解释。
\end{example}

在这个例子里，原环境可用的策略可以确定，跨环境重用的奖励却不能确定。恢复原策略、解释专家奖励和恢复可迁移奖励，是强度不同的结论，第三种还要支持新选择的排序。若新状态始终没有反馈，训练得到唯一参数解也可能只是正则化或模型结构替我们作了选择，而不是数据识别了唯一答案。

要尝试跨环境重用奖励，一种办法是分开任务计分与依赖原环境的价值信息。\term{对抗逆强化学习}（adversarial inverse reinforcement learning, AIRL）为此将奖励项和塑形项分别参数化，再用专家与策略转移训练判别器。例如采用
\begin{equation}
  \begin{aligned}
    D(s,a,s')&=\frac{\exp f(s,a,s')}
                         {\exp f(s,a,s')+\pi(a\mid s)},\\
    f(s,a,s')&=g(s,a)+\gamma h(s')-h(s).
  \end{aligned}
  \label{eq:rl-reward-airl}
\end{equation}
其中$g$表示待恢复的奖励项，$h$表示状态上的塑形评分，$f$组合两者；$D$仍取“来自专家”的概率方向。为说明实际更新，令本轮策略$\mu=\pi^{(k)}$，用$\bm\omega$参数化$g,h$，并在式~\eqref{eq:rl-reward-airl}的分母中令$\pi=\mu$，得到$D_{\bm\omega,k}$。专家与本轮策略转移按共同的折扣和吸收协议采样，分别标为1和0。保持$\mu$固定，最大化平衡来源样本的判别对数似然
\[
  \mathcal L_D(\bm\omega)
  =\mathbb E_{\mathrm E}\log D_{\bm\omega,k}(s,a,s')
    +\mathbb E_k\log[1-D_{\bm\omega,k}(s,a,s')] .
\]
两个期望分别对应专家和本轮策略的转移分布；经验更新用各来源批次的平均替代期望，再对$\bm\omega$作梯度上升。两种来源记录的分母都使用本轮$\mu(a\mid s)$，不是在专家记录中换成专家动作概率；在所用状态动作支持上须有$\mu(a\mid s)>0$。

策略更新时冻结学出的$g,h$，原论文采用判别器的对数几率作为奖励，而非直接使用$D$或只使用$f$。由式~\eqref{eq:rl-reward-airl}可得本轮固定的信号
\[
  \widetilde r_k(s,a,s')
  =\log D_{\bm\omega,k}-\log(1-D_{\bm\omega,k})
  =f_{\bm\omega}(s,a,s')-\log\mu(a\mid s).
\]
这一等式说明判别信号包含塑形奖励及温度1的动作熵项；使用该信号时不能再重复加同一熵奖励。用本轮轨迹计算它的回报及优势估计$\widehat A_t$后，可按
$\bm\theta^{(k+1)}=\bm\theta^{(k)}
+\eta\widehat{\mathbb E}_k[
\nabla_{\bm\theta}\log\pi_{\bm\theta}(a_t\mid s_t)\widehat A_t]$
作一次策略梯度更新，梯度在旧参数处计算，$\eta>0$为步长，样本平均按本轮折扣占用加权；多步复用可采用配套的PPO等代理更新。本轮将奖励和优势视为固定，下一轮才刷新策略、转移及判别器分母，不能把冻结$\mu$的代理目标说成任意新策略下的精确熵目标。

迁移时希望重用$g$，而不是整个$f$或上述判别信号。该解释有条件：原论文的恢复分析将一般形式进一步限制为状态奖励$g(s)$，并依赖真实奖励只依赖状态、确定性动态、可分解连通条件、理想最大因果熵专家及优化等假设，且仍有允许的常数歧义。这里的连通条件是：能由同一前驱状态一步到达的两个状态视为相连，再作传递闭包，要求所有状态相连；它使当前状态项与后继状态项能够分离，而不只是要求沿路径可达。一般状态动作奖励、随机动态或有限次优示范，不会仅因写出这个分解就获得无条件恢复保证\scite{rl5-fu2018airl}

这并不意味着所有势差一换环境就失效。若保留依赖实际转移的形式
\[
  R(s,a,s')+\gamma\Phi(s')-\Phi(s),
\]
并用新环境实际后继重新计算势差，满足原边界等条件时仍可相消。若先在旧环境对后继势取期望，保存成固定状态动作奖励
$R(s,a)+\gamma\mathbb E_{P_{\mathrm{old}}}\Phi(s')-\Phi(s)$，
再直接用于新动态，旧转移期望通常不再与新轨迹上的相邻势值相消。AIRL试图区分的正是与旧动态及专家行为纠缠的部分；可否迁移需检查实际奖励形式，不能将失败归于势差本身。

\begin{table*}[htbp]
  \centering\small
  \begin{tabularx}{\linewidth}{@{}>{\raggedright\arraybackslash}p{25mm}*{3}{>{\raggedright\arraybackslash}X}@{}}
    \toprule
    根因 & 可观察症状 & 不能仅凭什么判断 & 补充验证 \\
    \midrule
    代理定义遗漏 & 高分行为违反未编码要求 & 训练集评分准确率 & 要求到证据审计、针对遗漏的反例 \\
    目标信息不充分 & 多个相容奖励在新行为上分歧 & 奖励参数不唯一 & 有区分力的专家选择或偏好查询 \\
    反馈模型分布外误差 & 新策略行为评分与独立评价不符 & 旧分布拟合误差小 & 新行为重标、覆盖与校准检查 \\
    优化器利用漏洞 & 更新后反复产生某种虚假高分行为 & 单次异常或代理单独上涨 & 固定漏洞复现、版本对照、修补后再测 \\
    \bottomrule
  \end{tabularx}
  \caption{同一种行为偏差可能来自多个根因。奖励歧义本身不等于目标错误，唯一拟合也不能证明标准正确。}
  \label{tab:rl-reward-causes}
\end{table*}

表~\ref{tab:rl-reward-causes}将信息不足、标准遗漏、分布外误差和漏洞利用分开，后两者还可能相互强化。\term{奖励投机}（reward hacking）是策略利用代理缺陷取得高分，却违反原行为要求；\term{过度优化}（overoptimization）是继续加强代理优化时，独立目标可能转而恶化。更新改变行为分布，代理在新区域评分出错，优化又偏向虚假高分行为，就可能形成这种反馈。不过分布外误差未必已被利用，奖励不唯一也未必造成错误行为，仍要用具体反例和独立要求作判断。

% Outline: 5.4.4
\subsection{修订行为目标并约束策略更新}
\label{sec:rl-reward-revision}

发现高分配送有问题后，需要按根因选择修订对象。新增“内部完好”要求，是补充原目标；为同一完好标准补充标签说明，是修订操作性标注规则。重训评分器改进对既定标准的估计，修补允许空字符串通过的验证器则修正实现。四种改动影响不同，必须保留目标、标注、验证器和策略版本，不能统称为奖励模型升级。

修订以后若仍有估计误差，可以惩罚可估计的模型不确定性，用多个奖励模型的分歧定位缺口，或限制策略更新幅度。这些措施各有边界：模型可能共享盲点，不确定性估计也可能失准；保守更新减少进入未见区域的速度，同时可能压低原本可获得的收益。KL约束只控制相对参考策略的偏移，不能补回参考策略已遗漏的要求，平均KL小也不保证每个罕见状态都不变。

新增数据应回答已经发现的缺口。例如旧签收照片骗过代理时，可以补充照片与当前包裹身份的一致性检查，再收集不同伪造变体。模型刷新后重测原反例，并检查未参与修补的新变体，才能判断修订是否扩大了正确覆盖。比较时还要记录独立任务指标的变化、采用保守更新后减少的任务收益以及新增反馈成本，不能只报告旧训练集拟合改善。

图~\ref{fig:rl-reward-revision}显示不同修订如何返回检测：目标改动要重新明确要求，估计或实现改动要重新检验新行为。关闭一个反例不表示其他要求都已满足；风险、约束违反和运行时保护还需要第~\ref{chap:rl-trustworthy}章的工具。

\begin{figure*}[htbp]
  \centering
  % Revision is a diagnostic loop, not an additional feedback category.
\begin{tikzpicture}[
  x=1mm,y=1mm,font=\figurefont\fontsize{8}{10}\selectfont,
  box/.style={draw=BookRule,fill=BookPaper,text width=32mm,
    minimum height=12mm,align=center,inner sep=2mm,line width=.7pt},
  flow/.style={-{Stealth[length=2mm]},draw=BookMuted,line width=.8pt},
  lab/.style={fill=white,inner sep=1mm,font=\figurefont\fontsize{7}{9}\selectfont}
]
  \node[box,fill=FigureInput!10] (goal) at (20,36) {目标要求与训练代理};
  \node[box,fill=FigureOutput!10] (check) at (81,36) {独立检测};
  \node[box] (diagnose) at (143,36) {根因判断};
  \node[box,fill=FigureModel!10] (revision) at (143,0)
    {修订要求／模型／规则\\或限制策略更新};
  \node[box] (behavior) at (81,0) {更新后新行为};
  \node[box] (version) at (20,0) {版本对应与反例记录};
  \draw[flow] (goal) -- node[lab,above] {要求与检查对象} (check);
  \draw[flow] (check) -- node[lab,above] {具体反例} (diagnose);
  \draw[flow] (diagnose) -- node[lab,right] {待修补缺口} (revision);
  \draw[flow] (revision) -- node[lab,above] {新策略} (behavior);
  \draw[flow] (behavior) -- node[lab,left] {重新检测} (check);
  \draw[flow] (revision.south) -- ++(0,-12) -| node[lab,near start,below]
    {保留改变了什么} (version.south);
  \draw[flow] (version) -- node[lab,left] {修订后的约定} (goal);
\end{tikzpicture}

  \caption{目标与代理的诊断回路。检测传递具体反例，根因判断决定改变要求、修订估计还是限制更新；修订后的策略产生新行为，必须重新接受独立检测。}
  \label{fig:rl-reward-revision}
\end{figure*}

% Outline: 5.5
\section{三类奖励限制叠加时的反馈学习}
\label{sec:rl-reward-loop}

同一配送训练可能同时只有少量示范、部分路线偏好、过程记录和不完备的验证器。我们将这些材料接成一轮可复查的学习过程：明确本轮要求和反馈用途，更新有依据的模型或策略，再独立检验新行为，决定停止或补充信息。这里才有实际轮次的先后依赖；三类限制提供不同诊断依据，并不要求每轮依次消除它们。若反馈始终没有涉及某项要求，迭代也不能自行解除目标歧义。

% Outline: 5.5.1
\subsection{整合已有反馈形成当前反馈约束}
\label{sec:rl-reward-integration}

将配送示范、偏好和检查记录放在一起以前，需要知道它们各自在评价什么。一条可用记录至少包含任务输入、被评价行为、反馈值和形式、来源、采样策略及标准版本。轨迹评分还要记录步骤范围和终止原因，偏好还要保留两个候选及呈现顺序。这些字段使我们能够区分冲突来自行为条件、评分尺度还是标准变化。

这里用“反馈约束”概括记录对候选行为和评价模型提供的限制，但实施时必须分开三件事。行为要求规定什么算完成任务，例如交给正确收件人；训练目标规定怎样利用记录拟合模型，例如动作负对数似然和偏好比较损失。策略更新限制则规定本轮允许多大偏移，或排除哪些行为。带噪偏好通常进入损失，不是每条都必须满足的硬不等式；身份要求若作为硬约束，还要有可检查条件和可靠执行机制。

表~\ref{tab:rl-reward-records}按记录、约束对象和本轮用途整合配送反馈。同一数据可以支持多个模型，但动作标签、相对排序和正确性判断的语义不同，在说明聚合依据前不能直接相加成总分。

\begin{table*}[htbp]
  \centering\small
  \begin{tabularx}{\linewidth}{@{}*{5}{>{\raggedright\arraybackslash}X}@{}}
    \toprule
    反馈记录 & 约束对象 & 可靠范围 & 冲突处理 & 本轮用途 \\
    \midrule
    专家装载动作 & 条件动作分布 & 指令完整的已见状态 & 检查隐含指令与次优示范 & 行为克隆初始化或局部纠正 \\
    同包裹路线偏好 & 相对效用 & 同一偏好群体与标准 & 保留来源分层，不强行平均 & 拟合路线效用差 \\
    装载身份检查 & 步骤关系正确性 & 可读编号且检查器有效 & 与终点身份联合复核 & 提供局部信号 \\
    签收验证 & 最终身份条件 & 签收记录可信 & 排除旧记录与解析漏洞 & 终局反馈；另留独立样本 \\
    少量完好评分 & 包裹状态评价 & 评分尺度一致 & 复核尺度与开箱依据 & 校准反馈模型 \\
    独立开箱审查 & 未覆盖要求的证据 & 被抽查任务范围 & 不用于本轮拟合 & 检查代理遗漏 \\
    \bottomrule
  \end{tabularx}
  \caption{同一轮反馈记录的五项对应。独立证据应按样本和用途隔离，不能一边参与拟合一边称为留出检验。}
  \label{tab:rl-reward-records}
\end{table*}

假设专家走楼梯，偏好却选择电梯。先比较包裹是否易碎、当时电梯是否可用、评价者是否强调节能；只有条件和标准相同时，才需要进一步判断来源可靠性或专家是否次优。装载检查通过而终局身份失败，也可能来自中途换包，局部标签各自都正确，但跨步骤一致性被破坏。保留条件和粒度才能定位这些差异，平均正负分反而可能掩盖原因。

整合后，本轮要输出明确的要求、训练目标和更新限制。例如先要求身份正确，再在满足身份条件的行为中比较完好与能耗；用动作示范拟合策略，用路线偏好和完好评分分别拟合比较模型及评分模型。策略更新还可以采用参考KL惩罚或约束，但必须说明是哪一种。身份优先的文字约定并没有自动提供动作过滤器，硬性执行能力仍要验证。所有约定都应附上适用输入范围及尚未解决的分歧。

若分歧确实来自不能兼容的真实偏好，可以保留按偏好条件选择的策略或多个候选。必须聚合时，应解释分数单位、尺度、权重及权重代表的取舍；加权平均是一个新约定，不是争议已经被识别为唯一真实目标的证据。

% Outline: 5.5.2
\subsection{依据反馈约束联合更新模型与策略}
\label{sec:rl-reward-joint-update}

已有反馈及其用途确定后，要明确本轮改动哪些学习对象。对应记录可用于拟合奖励、效用或通过概率模型，归因机制再把轨迹结果转成局部信号，策略在本轮惩罚或约束下利用这些信号更新，并产生下一批行为。价值基线帮助降低方差或自举，不替代已经约定的行为评价标准。“联合更新”可以是交替拟合，也可以只改其中一项，不要求全部参数同时反向传播。图~\ref{fig:rl-reward-update}标明数据依赖，以及本轮冻结的标准和参考策略。

\begin{figure*}[htbp]
  \centering
  % Solid arrows carry data/signals; dashed boxes denote frozen objects.
\begin{tikzpicture}[
  x=1mm,y=1mm,font=\figurefont\fontsize{8}{10}\selectfont,
  box/.style={draw=BookRule,fill=BookPaper,text width=30mm,
    minimum height=11mm,align=center,inner sep=2mm,line width=.7pt},
  flow/.style={-{Stealth[length=2mm]},draw=BookMuted,line width=.8pt},
  lab/.style={fill=white,inner sep=.8mm,font=\figurefont\fontsize{7}{9}\selectfont}
]
  \node[box,fill=FigureInput!10] (records) at (18,28) {已整合反馈\\记录与版本};
  \node[box,fill=FigureModel!10] (model) at (80,28) {反馈模型\\奖励／效用／验证};
  \node[box] (attribution) at (143,28) {归因机制\\轨迹结果到局部信号};
  \node[box] (trajectory) at (18,-3) {新轨迹\\状态、动作、结果};
  \node[box,fill=FigureModel!10] (policy) at (80,-3) {策略\\$\pi^{(k+1)}$};
  \node[box] (signal) at (143,-3) {回报与优势\\价值基线可单独更新};
  \node[box,dashed,fill=white] (standard) at (80,56) {本轮冻结：目标标准};
  \node[box,dashed,fill=white] (ref) at (80,-31) {本轮冻结：参考策略};
  \draw[flow] (records) -- node[lab,above] {拟合数据} (model);
  \draw[flow] (model) -- node[lab,above] {评分} (attribution);
  \draw[flow] (attribution) -- (signal);
  \draw[flow] (signal) -- node[lab,above] {学习信号} (policy);
  \draw[flow] (policy) -- node[lab,above] {执行采样} (trajectory);
  \draw[flow] (trajectory.west) -- ++(-4,0) |- (records.west);
  \draw[flow] (trajectory.north) -- (18,10)
    -- node[lab,above] {轨迹结果} (111,10) |- (attribution.west);
  \draw[flow] (standard) -- node[lab,right] {语义约束} (model);
  \draw[flow] (ref) -- node[lab,right] {偏移约束} (policy);
  \draw[flow] (trajectory.south) -- ++(0,-12)
    node[below,align=center] {另送独立检验\\不回流本轮拟合};
\end{tikzpicture}

  \caption{一轮模型与策略更新的数据依赖。反馈记录约束反馈模型，轨迹结果经归因机制形成局部信号，策略产生下一批轨迹；虚线边界表示本轮冻结的目标标准和参考策略。检验数据不沿训练箭头进入拟合。}
  \label{fig:rl-reward-update}
\end{figure*}

在这条更新路径中，反馈来源和优化器是两个选择。RLHF与RLAIF可以采用相同策略优化器而使用不同评价者；迭代式偏好优化可以每轮补新偏好，再运行DPO；在线RLVR则对当前策略样本即时调用规则验证。解耦剪裁与动态采样策略优化（Decoupled Clip and Dynamic sAmpling Policy Optimization, DAPO）是其中一种具体更新方法。表~\ref{tab:rl-reward-loops}比较各轮更新和固定的对象，以及样本和反馈何时刷新，而不是将这些名称当成互斥算法类别。

\begin{table*}[htbp]
  \centering\small
  \begin{tabularx}{\linewidth}{@{}>{\raggedright\arraybackslash}p{23mm}*{3}{>{\raggedright\arraybackslash}X}@{}}
    \toprule
    路线 & 本轮更新对象 & 本轮固定对象 & 采样与反馈刷新 \\
    \midrule
    PPO式RLHF & 奖励模型、策略及价值可分阶段更新 & 标注标准、参考策略；策略阶段冻结奖励 & 当前策略采样，按预算补人类偏好 \\
    RLAIF & 反馈模型与策略，依实现可有价值模型 & AI评价模型、提示与标准的指定版本 & 新样本交由AI评价，独立抽查 \\
    迭代式偏好优化 & 偏好对上的策略参数 & 单轮偏好集与参考策略 & 外层采样和标注刷新，内层可离线DPO \\
    在线RLVR & 策略；基线依更新器而定 & 验证规则、答案或测试集 & 新策略候选逐批验证 \\
    DAPO & 无独立评论家的策略更新 & 当前验证规则、旧策略批次概率 & 动态采样补充混合得分组，按token聚合损失 \\
    \bottomrule
  \end{tabularx}
  \caption{反馈闭环而非互斥算法分类。DAPO是在线可验证反馈的一种具体实现，RLAIF也可以组合不同的策略优化器。}
  \label{tab:rl-reward-loops}
\end{table*}

表中的DAPO修改集中在剪裁、样本筛选和损失聚合。它在所讨论训练设置中使用规则结果奖励，去掉参考KL项，采用不同的上下剪裁阈值，并动态补充组内同时有成功和失败的输入。损失还从式~\eqref{eq:rl-reward-grpo}的回答内平均，改为对全部有效token平均。为比较长度权重，记剪裁项为$\ell_{i,t}$，两种聚合是
\begin{equation}
  \frac1m\sum_i\frac1{T_i}\sum_t\ell_{i,t},
  \qquad
  \frac1{\sum_iT_i}\sum_i\sum_t\ell_{i,t}.
  \label{eq:rl-reward-dapo-weight}
\end{equation}
前者赋予每个回答相同总权重，后者赋予每个token相同权重，使长回答在整批损失中占更大总权重。超长惩罚与截断处理也会改变反馈语义；论文称之为reward shaping，不意味着满足本章的潜势型不变条件。DAPO发表于NeurIPS 2025，但经验结果只对应报告中的模型、数学任务和预算，不能推出这些组件对所有任务都有收益\scite{rl5-yu2025dapo}

动态采样补混合得分组，可以减少零相对信号，但筛选也改变训练输入分布，被丢弃候选仍消耗采样和验证预算。过滤整组失败的困难题，没有解决这些题。提高剪裁上界只是改变代理目标梯度何时饱和，不是对动作概率施加硬上限。因而比较实现时，应同时记录样本选择、组内标准差分母、长度聚合方式和验证器版本。

\begin{example}[一轮可复查的更新记录]
  \label{ex:rl-reward-round}
  假设教学轮次$k=3$使用数据$D3$、身份与完好标准$S2$、反馈模型$M3$、旧策略$P3$和参考策略$P0$。$D3$有20条装载动作、12对偏好、8条完好评分，分别用于策略拟合、比较模型拟合和评分模型拟合，后两项组成$M3$。策略更新时冻结$S2,M3,P0$，独立审查记录不加入$D3$。

  再取两条等长四步轨迹，以本轮约定的综合代理$\widetilde R$给终局记1和0分，中间奖励为零，$\gamma=0.9$，于是起点回报为$0.729$和0。冻结基线为$0.3645$，减去基线得到样本优势$0.3645$和$-0.3645$。只更新起点二动作的logit，其他动作规则固定；初始logit为0，两条轨迹各来自一个起点动作，样本平均梯度为$0.18225$。令步长为$0.2$、KL系数为$0.1$，且当前策略恰与参考相同，精确KL梯度此时为零。一次未剪裁更新得到logit $0.03645$，正优势动作概率约为$0.5091$，作为候选$P4$。

  本轮记录目标为“轨迹代理回报的PPO式改进减参考KL”，并分别预留40个访问状态动作标签、20对比较、100次验证和10次人工审查的预算上限。更新后要分开检验$M3$是否更符合反馈和$P4$是否改善实际行为。这些计算是可复查的教学记录，不是实测性能。
\end{example}

% Outline: 5.5.3
\subsection{用更新后行为检验反馈约束的有效性}
\label{sec:rl-reward-loop-check}

候选$P4$产生新行为后，要按本轮开始前的标准检查，不能只重算$P3$的旧日志分数。三类限制分别对应三个问题：新状态是否超出反馈覆盖，局部信号是否仍正确联系关键动作与结果，代理回报增加时独立行为指标是否改善。表~\ref{tab:rl-reward-checks}分列相应证据和失效判据，使一种检查通过不能替代其他检查。

\begin{table*}[htbp]
  \centering\small
  \begin{tabularx}{\linewidth}{@{}>{\raggedright\arraybackslash}p{20mm}*{4}{>{\raggedright\arraybackslash}X}@{}}
    \toprule
    检查 & 输入行为 & 独立证据 & 失效判据 & 仍未覆盖范围 \\
    \midrule
    观测与覆盖 & 新策略访问状态、长短轨迹 & 新状态专家标签、留出来源评分 & 无支持区域仍被赋予高置信评价 & 未采到的输入和动作 \\
    局部归因 & 关键步骤与相关标记 & 受控扰动、过程复核、终局结果 & 无关标记改变信号，关键错误不受罚 & 不可干预步骤的因果贡献 \\
    目标匹配 & 代理高分的新行为 & 独立身份、完好和成本检查 & 违反已约定要求或独立指标退化 & 未形式化与难观测要求 \\
    来源一致性 & 同行为的多来源判断 & 第三方执行或事实核对 & 分歧无法由条件差异解释 & 一致来源的共同盲点 \\
    \bottomrule
  \end{tabularx}
  \caption{检验必须说明未覆盖范围。来源一致性可以暴露冲突，却不能替代外部事实证据。}
  \label{tab:rl-reward-checks}
\end{table*}

检验结果应区分证据支持、证据不足和明确失效。若$P4$在预定短路线、身份可读范围内完成全部检查，估计精度也达到事先要求，可以支持当前反馈在这个范围有效。若只测三次且均成功，仍未覆盖必要长路线或满足精度要求，就只能说暂未发现失败。若已有代理高分却交错收件人的行为，便得到明确反例；平均成功率高不能抹去这个失效。

诊断要返回对应缺口：新状态无标签就补访问覆盖，关键步骤局部信号错误就修订归因，完好要求未编码就修订标准或代理。这样可以针对已有证据改变相关对象，不必一次失败就重训全部模型。反例若进入下一轮修订，就成为开发数据，后续不能再将它单独称为独立留出检验。

% Outline: 5.5.4
\subsection{依据检验结果与反馈预算决定停止或继续}
\label{sec:rl-reward-stopping}

即使新行为全部通过已做检查，也需要判断检查是否达到了事先约定的要求。训练前应固定独立指标的接受标准、必要覆盖、统计精度或不确定性要求，以及各类预算上限，不能看到结果后再放宽门槛。有限测试可以支持限定范围的接受，不为未测范围提供“没有失败”的保证。

\begin{table*}[htbp]
  \centering\small
  \begin{tabularx}{\linewidth}{@{}>{\raggedright\arraybackslash}p{23mm}*{3}{>{\raggedright\arraybackslash}X}@{}}
    \toprule
    证据状态 & 允许的后续动作 & 必须保留的结论范围 & 所需信息 \\
    \midrule
    达到预设充分性与接受标准 & 停止训练，保留复核入口 & 只接受已测任务、标准和策略版本 & 已完成覆盖、精度及成本记录 \\
    未发现失败但证据不足 & 有预算则定向查询；否则停止或收窄范围 & 未证实原范围满足要求；预算停止不等于成功 & 缺少的场景、样本或独立来源 \\
    已发现明确失效 & 暂停对该范围的接受，修订后再测；或终止 & 记录反例与不适用范围 & 根因、修补方案、验证与查询成本 \\
    \bottomrule
  \end{tabularx}
  \caption{停止训练是一种行动，接受行为是一个证据结论，二者不能混同。}
  \label{tab:rl-reward-stopping}
\end{table*}

表~\ref{tab:rl-reward-stopping}区分停止行动和接受结论。满足预设要求可以停止，预算耗尽也可以停止，但后者仍可能留有关键缺口。若获准查询只能看终点照片，不能观察包裹内部，重复评价也无法识别完好要求；这时应收窄结论或改变信息接口，而不是将次数当作充分性。

有剩余预算时，还要说明下一批反馈可能改善哪项判断：当前状态覆盖、候选目标排序、过程贡献或验证规则。可以估计查询减少多少候选模型分歧、是否可能改变决策，再与人工时间、采样和执行成本比较；没有校准数据时，这些只是启发式依据，不是有数值保证的信息价值。更换目标、反馈标准或明显改变策略访问分布后，原停止证据所对应的对象也变了，需要重新检查。

\begin{algorithm}[反馈学习的继续与停止闭环]
  \label{alg:rl-reward-loop}
  \begin{algorithmic}[1]
    \Require 目标标准、独立接受条件、必要覆盖、分项预算、初始反馈与策略
    \Ensure 候选策略、结论范围、停止原因和未解决缺口
    \While{尚未记录停止原因}
      \State 整合兼容的反馈，固定本轮标准、数据范围及未解决分歧
      \State 只更新有相应证据支持的反馈模型、归因机制或策略
      \State 以新行为接受独立检查，分类为充分、证据不足或明确失效
      \If{覆盖与精度充分，且满足预设接受条件}
        \State 记录“在已检验范围满足要求”，停止本轮学习
      \ElsIf{预算不足，或允许查询不能区分关键目标解释}
        \State 记录“预算停止”或“识别受限”，收窄结论并停止
      \Else
        \State 按观测、归因或目标缺口选查询，必要时先修订目标标准
        \State 获取新反馈，记录成本、生成策略与标准版本，返回整合
      \EndIf
    \EndWhile
  \end{algorithmic}
\end{algorithm}

算法~\ref{alg:rl-reward-loop}只更新有相应证据支持的对象，不要求每轮全部模型变化。明确失效可能先需要修订要求，而不是马上增加样本；若获准查询无法补足缺口，也可以停止优化错误代理，并保留未解决结论。这个规则控制反馈学习是否继续，不替代系统部署审批或运行流程。

% Outline: 5.5.5
\subsection{针对反馈缺口获取新反馈}
\label{sec:rl-reward-query}

决定继续后，还需把未决问题转成可执行查询。例如错误走廊缺恢复动作，就查询该状态上的专家标签；两条新路线排序不定，就比较候选奖励意见不同的路线。表~\ref{tab:rl-reward-queries}按这类缺口列出查询能排除的解释，同时保留查询之后仍未回答的部分。

\begin{table*}[htbp]
  \centering\small
  \begin{tabularx}{\linewidth}{@{}*{4}{>{\raggedright\arraybackslash}X}@{}}
    \toprule
    具体缺口 & 查询方式 & 可以排除的解释 & 仍不能回答什么 \\
    \midrule
    偏离后不会恢复 & 当前访问状态上的专家动作 & 这些状态上的错误动作候选 & 唯一真实奖励与未查询状态 \\
    新路径排序不确定 & 比较奖励候选意见不同的路径 & 对该比较给出相反排序的目标 & 其他输入或群体的偏好 \\
    关键过程贡献不清 & 步骤标注、受控检查与结果对照 & 局部无效却被奖励的归因 & 未控制混杂下的完整因果结构 \\
    验证标准遗漏 & 新结果检查、规则边界与对抗案例 & 会接受已知错误行为的验证器 & 未被形式化的全部要求 \\
    来源共同偏差 & 独立来源和实际执行审查 & 仅依赖共同代理的解释 & 未审查行为的真实性 \\
    \bottomrule
  \end{tabularx}
  \caption{查询应增加能区分候选解释的信息，而不只是重复已经一致的反馈。}
  \label{tab:rl-reward-queries}
\end{table*}

查询的选择依据也不同。DAgger查询当前访问状态的专家动作，主动偏好查询选择可能改变候选排序的比较；不确定性采样选择反馈模型分歧大的行为，在线验证检查当前策略产物，对抗收集则寻找代理漏洞。但模型分歧不必等于真实目标信息量：共同盲点可能让模型一致自信，模糊措辞却可能造成很大分歧而不改变行为要求。保留独立随机或分层抽查，可以检查主动筛选是否遗漏重要区域。

沿示例~\ref{ex:rl-reward-round}，假设$P4$主要缺恢复动作，可以查询6个实际访问状态的动作；为新路线再补4对偏好、24次验证和2次人工审查。还要分别记录取得候选所用轨迹数及环境步数。这些教学计数单位不同，相加成“36条反馈”不能说明成本；每批记录应附带生成策略$P4$、标准$S2$及采样筛选协议。

新反馈返回第~\ref{sec:rl-reward-integration}节时，先核对与旧标准的兼容性。沿用$S2$的记录可聚合；若$S3$新增内部破损要求，旧终点照片标签只支持原先被检查的条件，需视用途重标或分开使用，不能直接当作$S3$的完整正例。若全部获准查询只比较速度，就仍无法区分看重完好和看重能耗的目标解释。此时应报告识别边界，闭环轮数没有提供缺失的信息。

% Outline: 5.6
\section{本章小结}
\label{sec:rl-reward-summary}

回到配送中的反馈缺口，替代记录首先决定我们能学到什么。示范约束动作或长期占用，评分和偏好约束评价函数，验证只判断明示条件；它们不相互等价，也不必都先还原为唯一奖励。BC在已有专家分布上拟合，DAgger补当前访问状态；奖励模型加策略优化需要评分新行为，标准DPO拟合固定偏好对。GRPO怎样比较组内候选，则与分数是否来自规则是两个选择。

取得结果以后，还要判断它怎样支持早期动作更新。结果回溯重新记账，过程监督增加证据，奖励变换和辅助目标改变信号或训练位置，各有成本与保持条件。潜势塑形在有界势、共同初始分布及正确终端处理下，靠相邻势项相消保持标准期望回报排序；这不证明局部信号具有因果意义，也不保证风险排序或学习速度不变。

最后，训练信号是否表达行为要求，必须由独立证据检查。对全部策略保持排序的奖励变换只改变表示；只解释同一批反馈的候选奖励却可能在未比较行为上分歧。拟合成功不证明标准完整，奖励歧义也不必妨碍选择当前范围内的可用策略。反馈学习能否继续改进，取决于每轮是否明确固定标准、更新对象、检验证据和新增信息；停止时还要区分满足限定要求、预算耗尽和识别受限。更密集反馈、更大模型或更多轮次，都不能代替这些判断。
